\documentclass[output=paper,colorlinks,citecolor=brown]{langscibook}
\ChapterDOI{10.5281/zenodo.21237005}
\author{David Sasseville\orcid{}\affiliation{École Pratique des Haute Études, Paris}}
%\ORCIDs{}

\title{The relationship between \emph{u}-adjectives and \emph{nu}-causatives in Anatolian}

\abstract{The significant number of Hittite causative/factitive formations in -\textit{nu}- attested next to \textit{u}-adjectives has always been striking and led together with comparative evidence from Old Indic and Ancient Greek to a theory in the 70s that the process occurs through \textit{n}-infixation between the root and the adjectival suffix -\textit{u}-, e.g., Hitt.\ \textit{tēp-u}- ‘small’ → Hitt.\ \textit{tēp-n-u-\textsuperscript{mi}} ‘to make small, diminish’ (matching OInd.\ *\textit{dabh-u-} (< PIE *\textit{dʰebʰ-u}-, replaced by OInd.\ \textit{dabh-rá-} ‘small’) → OInd.\ \textit{dabh-n-ó-ti} ‘to harm, deceive’), Hitt.\ \textit{park-u-} ‘high’ (< PIE *\textit{bʰerg̑ʰ-u}-, beside Hitt.\ \textit{park-taru} ‘must become high’) → Hitt.\ \textit{parg-n-u-\textsuperscript{mi}} ‘to make high’, provoking at a certain point a reanalysis of the morpheme  -\textit{n}- to -\textit{nu}-. Although this theory never gained much popularity, recent studies on the Luwian language have brought forth more data that have yet to be examined in depth vis-à-vis the morphological process involved, e.g., HLuw.\ \textit{wa/i-su}- /wasu-/ ‘good’ (< PIE *\textit{h\textsubscript{1}u̯es-u}-, beside HLuw.\ \textit{wa/i\nobreakdash-s}˚ /wass-\textsuperscript{ti}/ ‘to be good, pleasing’) → HLuw.\ \textit{u-sa-nú}- /usnuwa-\textsuperscript{i}/ ‘to make good, benefit’, and, thus, a re-examination of the Anatolian material is warranted. In addition to retracing the development of the denominal morpheme -\textit{nu}-, its extension to other types of nominal stems, notably to the \textit{i}-stems (e.g., Hitt.\ \textit{ḫarki}- ‘white’ (< PIE *\textit{h\textsubscript{2}erg̑-i}-) → Hitt.\ \textit{ḫarg(a)nu}- ‘to make white’), as well as the motivation behind it is  addressed. Furthermore, the phenomenon of \textit{u}-extended roots that has been considered to be intertwined with the development of the \textit{nu}-suffix is  explored and the discussion is expanded with a new set of Anatolian data (cf.\ HLuw.\ \textit{wa/i\textsubscript{5}-su}- /wasu-/ ‘to be good, successful’).}


\IfFileExists{../localcommands.tex}{
   \addbibresource{../localbibliography.bib}
   \usepackage{tabularx,multicol}
\usepackage{url}
\urlstyle{same}

 
\usepackage{langsci-optional}
\usepackage{langsci-lgr}
\usepackage{langsci-gb4e}
\usepackage[linguistics,edges]{forest}
\usepackage{tikz-qtree}
\usetikzlibrary{positioning}
\usetikzlibrary{patterns.meta}
\usepackage{multirow}
\usepackage{pgfplots}
\usepackage{setspace}
\usepackage{amsmath,stackengine}
% \usepackage{fontspec}
% \usepackage{tipa} % TIPA is banned
\usepackage{langsci-textipa}
\usepackage{langsci-branding}
\usepackage{longtable}
\usepackage{tabto}

\usepackage{enumitem}

%\usepackage{lscape} %Querformat chapter 4

   % \newcommand*{\orcid}{}


\makeatletter
\let\thetitle\@title
\let\theauthor\@author
\makeatother

\newcommand{\togglepaper}[1][0]{
   \bibliography{../localbibliography}
   \papernote{\scriptsize\normalfont
     \theauthor.
     \titleTemp.
     To appear in:
    Laura Grestenberger, Viktoria Reiter \& Melanie Malzahn (eds.).
    \emph{Deadjectival verb formation in Indo-European and beyond}.
     Berlin: Language Science Press. [preliminary page numbering]
   }
   \pagenumbering{roman}
   \setcounter{chapter}{#1}
   \addtocounter{chapter}{-1}
}

\newbool{bookcompile}
\booltrue{bookcompile}
\newcommand{\bookorchapter}[2]{\ifbool{bookcompile}{#1}{#2}}

%Sonderzeichen
%idg palatales k
\newcommand{\kinvbreve}{k\raisebox{0.6ex}{\hspace{-0.05em}\char"0311}}

\newcommand{\jac}{%
  \textit{\char"0237}% italic dotless j (ȷ)
  \hspace{-0.01em}% fine-tune horizontal spacing
  \raisebox{0.01ex}{\char"0301}% raised combining acute
}

%Baltic circumflex l
\newcommand\ltilded{%
\hspace{-0.2em}%
  \stackon[-1.5pt]{l}{\smash{\kern0.2em\~{}}}%
  \hspace{-0.2em}%
}

%Slav. wedge+ double acute
\newcommand{\HighH}[1]{%
    \stackon[-1.2ex]{#1}{\kern0.1em\H{}\kern-0.1em}%
}

%Lith dot tilde VICKY diese b.-sl. Sonderzeichen BRINGEN MICH UM!!
\newcommand{\tildedot}[1]{%
    \ensurestackMath{%
        \stackon[-0.27ex]{% Tilde layer 
            \stackon[-0.15ex]{% Dot layer 
                #1%
            }{\hspace{0.1em}\cdot\hspace{-0.13em}}%
        }{\hspace{0.1em}\text{\~{}}\hspace{-0.13em}%
    }%
}}

%Lith dot acute
\newcommand{\dotacute}[1]{%
    \ensurestackMath{%
        \stackon[-1.1ex]{% acute layer 
            \stackon[-0.15ex]{% Dot layer 
                #1%
            }{\hspace{0.1em}\cdot\hspace{-0.12em}}%
        }{\hspace{0.1em}\text{\'{}}\hspace{-0.12em}%
    }%
}}



% \setlength{\emergencystretch}{2pt} 

 \pgfplotsset{compat=1.18}


\defcitealias{AbB_1}{AbB 1}
\defcitealias{AbB_7}{AbB 7}
\defcitealias{ABL_1}{ABL 1}
\defcitealias{Vries1977}{AEW}
\defcitealias{ALEW}{ALEW}
\defcitealias{AnGr}{AnGr}
\defcitealias{ARM_1}{ARM 1}
\defcitealias{BT}{BT}
\defcitealias{black_concise_2000}{CDA}
\defcitealias{CHD}{CHD}
\defcitealias{CT28}{CT 28}
\defcitealias{CT34}{CT 34}
\defcitealias{CT39}{CT 39}
\defcitealias{Cleasby1874}{CV}
\defcitealias{Beekes2010}{EDG}
\defcitealias{eDiAna}{eDiAna}
\defcitealias{eDIL}{eDIL}
\defcitealias{EDL}{EDL}
\defcitealias{Kroonen2013}{EDPG}
\defcitealias{ESJS}{ESJS}
\defcitealias{EVP}{EVP}
\defcitealias{EWA1988}{EWA}
\defcitealias{EWAI}{EWAia I}
\defcitealias{EWAII}{EWAia II}
\defcitealias{EWGP}{EWGP}
\defcitealias{GDLI1961}{GDLI}
\defcitealias{GPC}{GPC}
\defcitealias{GRADIT1999}{GRADIT}
\defcitealias{HW}{HW}
\defcitealias{IEW}{IEW}
\defcitealias{KAH_2}{KAH 2}
\defcitealias{KAR_1}{KAR 1}
\defcitealias{KAR_2}{KAR 2}
\defcitealias{KEWA}{KEWA}
\defcitealias{Labat_TDP}{Labat TDP}
\defcitealias{lambert_BWL}{Lambert BWL}
\defcitealias{LEIA}{LEIA}
\defcitealias{Rix2001}{LIV\textsuperscript{2}}
\defcitealias{LIVaddenda}{LIV\textsuperscript{2} {A}ddenda}
\defcitealias{LKG1971}{LKG}
\defcitealias{LKZ}{LKŽ}
\defcitealias{LP}{LP}
\defcitealias{MhdGr}{MhdGr}
\defcitealias{MLLVG1959}{MLLVG}
\defcitealias{NIL}{NIL}
\defcitealias{OgnS}{OgnS}
\defcitealias{ONP}{ONP}
\defcitealias{RIMA_2}{RIMA 2}
\defcitealias{TIM_2}{TIM 2}
\defcitealias{WAP}{WAP}
\defcitealias{YOS_2}{YOS 2}






\newindex{wan}{wdx}{wnd}{Form index}
\newcommand{\iw}[2]{\index[wan]{#1!{#2}\xspace}}
\newcommand{\iwi}[3][]{\index[wan]{#2!{#3}\xspace#1}{#3}}
\newcommand{\iwit}[3][]{\rule{0pt}{0pt}#1\index[wan]{#2!{#3}}{#3}}


\newindex{van}{vdx}{vnd}{Form index}
\newindex{ban}{bdx}{bnd}{Book index}

% \newcommand{\iw}[3][]{\index[wan]{#2!\textit{#3}}}


% \providecommand{\iwi}[3][]{\iw{#3}\textit{#3}}


\newcommand{\indexnote}[1]{{\footnotesize\sffamily\raggedright\normalfont\color{gray}#1}}

   %% hyphenation points for line breaks
%% Normally, automatic hyphenation in LaTeX is very good
%% If a word is mis-hyphenated, add it to this file
%%
%% add information to TeX file before \begin{document} with:
%% %% hyphenation points for line breaks
%% Normally, automatic hyphenation in LaTeX is very good
%% If a word is mis-hyphenated, add it to this file
%%
%% add information to TeX file before \begin{document} with:
%% %% hyphenation points for line breaks
%% Normally, automatic hyphenation in LaTeX is very good
%% If a word is mis-hyphenated, add it to this file
%%
%% add information to TeX file before \begin{document} with:
%% \include{localhyphenation}
\hyphenation{
    par-a-digm non-de-adj-ect-iv-al -pre-sent -pre-sents for-ma-tion-al in-do-ger-ma-ni-sche in-do-ger-ma-ni-schen Rei-chert Rei-ter dia-chrony ety-mo-lo-gi-sches Ost-ger-ma-ni-schen ger-ma-ni-schen alt-in-di-schen mor-pho-lo-gi-scher Alt-indo-ari-sch-en lexico-graphy Schrij-ver
}
% aktuell in Bibliographie: stud-ies, histor-ical Soci-età
\hyphenation{
    par-a-digm non-de-adj-ect-iv-al -pre-sent -pre-sents for-ma-tion-al in-do-ger-ma-ni-sche in-do-ger-ma-ni-schen Rei-chert Rei-ter dia-chrony ety-mo-lo-gi-sches Ost-ger-ma-ni-schen ger-ma-ni-schen alt-in-di-schen mor-pho-lo-gi-scher Alt-indo-ari-sch-en lexico-graphy Schrij-ver
}
% aktuell in Bibliographie: stud-ies, histor-ical Soci-età
\hyphenation{
    par-a-digm non-de-adj-ect-iv-al -pre-sent -pre-sents for-ma-tion-al in-do-ger-ma-ni-sche in-do-ger-ma-ni-schen Rei-chert Rei-ter dia-chrony ety-mo-lo-gi-sches Ost-ger-ma-ni-schen ger-ma-ni-schen alt-in-di-schen mor-pho-lo-gi-scher Alt-indo-ari-sch-en lexico-graphy Schrij-ver
}
% aktuell in Bibliographie: stud-ies, histor-ical Soci-età
   \boolfalse{bookcompile}
   \togglepaper[4]%%chapternumber
}{}

\begin{document}
\SetupAffiliations{mark style=none}
\maketitle



\section{Introduction}


The numerous examples of Hittite denominal\is{denominal!verb}  \isi{factitive}s in \iwi{Hitt}{-\textit{nu}-} \is{nasal present} had already been noticed in the 30s, when  comparative linguistic work on Hittite was thriving. Indo-Europeanists were convinced that the suffix \iwi{Hitt}{-\textit{nu}-} was primarily deverbal\is{deverbal/deverbative!verb} and that the denominal function \is{denominal!verb} had been secondarily acquired by Anatolian.\ In his doctoral thesis submitted in 1973, Harold Koch scrutinized the data more than it had been done before and cogently defended the view that Hittite \isi{factitive}s in \iwi{Hitt}{-\textit{nu}-} were originally derived from \textit{u}-adjectives\iw{Hitt}{-\textit{u}-} via \textit{n}-infixation.\ Further, he explained that the co-existence of root verbs, \textit{u}-adjectives and \textit{nu}-\isi{factitive}s all based on the same root would have been the starting point for the replacement of the isolated infixation process with the more common \isi{suffixation} one.\footnote{\citealt[25--26, 34--36]{Koch1973}. \citet[165--166 with fn.\ 72]{Oettinger1979} points out that the infixation process should be older than the apparent separation between nominal and verbal roots within Proto-Indo-European, adducing PIE \iwi{PIE}{*\textit{u̯ed}-} ‘water’ → \iwi{PIE}{*\textit{u-né-d-}/\textit{und}-´} ‘to water’ (continued in Ved.\ \iwi{Sanskrit}{\textit{unátti}}) as a case in point; cf.\ \citetalias[  712, fn.\ 43]{NIL}. However, the nature of Proto-Indo-European roots and the relationship between diverse root formations is a topic that will require more attention in future research.} Thus, by reanalysing\is{reanalysis} the derivational base, the deverbal\is{deverbal/deverbative!verb} \isi{causative} suffix would have been created.\footnote{This approach slightly differs from the theory defended by \citet[60--61]{Strunk1967}, who maintains that the suffix \iwi{Hitt}{-\textit{nu}-} originates in the \isi{reanalysis} of the infix to a root ending in -\textit{u̯} such as PIE \iwi{PIE}{*\textit{\kinvbreve leu̯}-} ‘to hear’ → \iwi{PIE}{*\textit{\kinvbreve l̥-né-u̯-/\kinvbreve l̥-nu}-´} ‘to cause to hear’. However, since this putative proto-type is rare (see also \citealt{VillanuevaSvenssonforth}), Strunk had to resort to \textit{u}-extended roots in order to provide further support to this theory (the role of \textit{u}-extensions for this topic is addressed in detail in the last part of this paper). The view presented here also differs from the more recent theory by \citet{Steer2015}, who argues that the \textit{nu}-suffix\is{nasal present} is primary and the \textit{n}-infix secondary. Although his approach is morphologically acceptable, it goes against the evidence of the daughter languages, in which the \textit{n}-infix is older than its various reanalyses\is{reanalysis} as a suffix.} This theory can be illustrated with the following examples, i.e., first the proto-type and second the \isi{reanalysis}:


\begin{enumerate}
   \item  {CVC\iwi{Hitt}{-\textit{u}-} → CVC-\textit{n-u}- (e.g., Hitt.\ \iwi{Hitt}{\textit{ḫuišu}-} ‘alive, raw’ → \textit{ḫuiš-n-u-\textsuperscript{mi}}\iw{Hitt}{\textit{ḫuišnu}-} ‘to keep alive’)}
\item CVC- → (\isi{reanalysis}) CVC\iwi{Hitt}{-\textit{nu}-} (e.g., Hitt.\ \textit{ḫuiš-\textsuperscript{mi}}\iw{Hitt}{\textit{ḫuiš}-} ‘to live’ → \textit{ḫuiš-nu-\textsuperscript{mi}}\iw{Hitt}{\textit{ḫuišnu}-} ‘to keep alive, let live’)
\end{enumerate}

It is noticeable that the formation of Hittite \textit{ḫuišnu-\textsuperscript{mi}}\iw{Hitt}{\textit{ḫuišnu}-} is synchronically ambiguous, as one could even switch the example of the proto-type with the one of the \isi{reanalysis}, arguing for the reverse direction in the development and making the \textit{n}-infixation younger than the \iw{Hitt}{-\textit{nu}-}\textit{nu}-\isi{suffixation}.\footnote{Thus \citealt{Steer2015}.} \citet[34--36]{Koch1973},   however, insists that the \textit{n}-infixation is older by pointing out that an \textit{n}-infixed verb like Hitt.\ \textit{ḫarni(n)k-\textsuperscript{mi}}\iw{Hitt}{\textit{ḫarni(n)k}-} ‘to destroy’ is later replaced by the more productive\is{productivity} suffixed\is{suffixation} formation Hitt.\ \textit{ḫarknu-\textsuperscript{mi}}\iw{Hitt}{\textit{ḫarknu}-} ‘to destroy’.\footnote{Thus also \citealt[192]{Oettinger1979}, \citealt[281--282]{Shatskov2020}.} The recently discovered Luwian verb \textit{tarnuwa-\textsuperscript{i}}\iw{CLuw}{\textit{tarnuwa}-} ‘to release’ vis-à-vis Hitt.\ \textit{tarna-/tarn-\textsuperscript{ḫi}}\iw{Hitt}{\textit{tarna-/tarn}-} ‘to release’ is another example of such a replacement (cf.\ \citealt[109--110]{Melchert2022}), comparable to the relationship between Ved.\ \iwi{Sanskrit}{\textit{str̥ṇā́ti}} ‘spreads’ and Gk.\ στόρνυμι\iw{AG}{στόρνῡμι} ‘id.’ or Ved.\ \iwi{Sanskrit}{\textit{yunákti}} ‘yokes’ and Gk.\ ζεύγνυμι\iw{AG}{ζεύγνῡμι} ‘id.’, where the original \textit{n}-infixed verb has been replaced in Ancient Greek by a \textit{nu}-suffixed\is{suffixation} verb (see \citealt[308--309]{Rasmussen1999} and \citealt{VillanuevaSvenssonforth} for a summary of these phenomena). The radical \textit{n}-infixation in the Anatolian branch is indeed an archaism that can still be observed in Hittite, but almost vanished in the other 2\textsuperscript{nd} millennium BC languages like Luwian and Palaic. In these other two languages, it appears to have been replaced by \isi{suffixation} after having been reanalyzed\is{reanalysis} with a root-final laryngeal before combining with another morpheme; the result being \textit{\nobreakdash-i\nobreakdash-na-\textsuperscript{i}}\iw{CLuw}{-\textit{i-na}-} \iw{HLuw}{-\textit{i-na}-} \iw{Palaic}{-\textit{i-na}-} (consisting of verbal \iwi{PIE}{*-\textit{i̯e/o}-} \is{*-\textit{i̯e/o}-present} or nominal \iwi{PIE}{*-\textit{i}-}, followed by either \iwi{PIE}{*-\textit{néh\textsubscript{2}-/-nh\textsubscript{2}\nobreakdash-}´} or \iwi{PIE}{*-\textit{néh\textsubscript{3}-/-nh\textsubscript{3}}-´}).\footnote{\citealt[540f., 548]{Sasseville2021}.} On the other hand, the \textit{nu}-suffix experienced an increasing \isi{productivity} in Hittite and Luwian, as will be shown in more detail in the next section.

In order to prove that the \textit{n}-infixation among \textit{u}-adjectives is inherited, \citet{Koch1973} cites parallels in Old Indic and Ancient Greek. The Old Indic example for the proto-type that has been cited in the literature for almost a century and was convincingly defended by \citet[164--165]{Oettinger1979} and \citet{Koch1980} is the following one:\footnote{\citealt[69]{GotzePedersen1934}, \citealt[231]{Kammenhuber1969}, \citetalias[  694--696]{EWAI}; followed by \citealt[164--165]{Oettinger1979}, \citealt[27]{Tucker1981}, \citealt[385--386]{Praust2004}, \citealt[208--210]{Steer2015}, \citealt[2162]{LundquistYates2018}, \citealt[157]{Oettinger2022} (adding PIE \iwi{PIE}{*\textit{méi̯-u}-} ‘few’ → \iwi{PIE}{*\textit{mi-né-u̯-ti}} ‘to make few’). The first objection raised by \citet[277]{Shatskov2020} that Hitt.\ \textit{tepnu-\textsuperscript{mi}}\iw{Hitt}{\textit{tepnu}-} cannot be old because it is not attested in Old Hittite is invalid. The “Old Hittite” corpus is very small and highly fragmentary, which makes Shatskov’s claim \emph{ex silentio}. Shatskov also argues against the word equation due to the, in his words, “quite different” meanings, although \citet[225--232]{Koch1980}  has explained in detail the semantic development of ‘to belittle’ → ‘to deceive’ and provided sufficient contextual parallels between the Hittite and the Indo-Iranian cognate (e.g., the gods as patient). It is also interesting to notice that the Hittite \isi{factitive} \iwi{Hitt}{\textit{tepnu}-} (just like OInd.\ \iwi{Sanskrit}{\textit{dabhnóti}}) shows the restored full grade of the \textit{u}-adjective. Given that the expected zero-grade of the root would have yielded an unstable consonant cluster, i.e., Hitt.\ [tbnu-], it may have facilitated a restoration of the full grade in analogy to the \textit{u}-adjective in the prehistory of each language. Only Old Avestan \iwi{Avestan}{\textit{dǝbǝnaotā}} preserves the original zero-grade with the complex consonant cluster \citep[164]{Oettinger1979}.} Hitt.\ \textit{tēpu-/-aw-}\iw{Hitt}{\textit{tēpu}-} ‘small’ → \textit{tepnu-\textsuperscript{mi}}\iw{Hitt}{\textit{tepnu}-}  ‘to make small, belittle’ matching OInd.\ *\textit{dabh-u-}\iw{Sanskrit}{\textit{dabhu}-} (replaced by OInd.\ \textit{dabh-rá-}\iw{Sanskrit}{\textit{dabhrá}-} ‘small’) → OInd.\ \textit{dabh-n-ó-ti}\iw{Sanskrit}{\textit{dabhnóti}} ‘to harm, deceive’, both leading to the PIE reconstruction \iwi{PIE}{*\textit{dʰ(é)bʰ-u-/-éu̯}-} ‘small’ → \iwi{PIE}{*\textit{dʰ(e)bʰ-n-éu̯-/-u}-´} ‘to make small’.

Whereas the denominal process\is{denominal!verb} has been given up in Old Indic, \citet{Koch1978} and \citet[22--28]{Tucker1981}  both argued independently that the process is continued in Ancient Greek \isi{factitive}s in -ύνω,\iw{AG}{-ῡ́νω} e.g., \iwi{AG}{βαθύς} ‘deep’ → βαθύνω\iw{AG}{βαθῡ́νω} ‘to deepen’, \iwi{AG}{θρασύς} ‘bold, courageous’ → \iwi{AG}{θαρσύνω} ‘to embolden’, \iwi{AG}{εὐρύς} ‘wide’ → \iwi{AG}{εὐρύνω} ‘to widen’ and \iwi{AG}{*αἰσχύς} (replaced by αἰσχ-ρό-ς\iw{AG}{αἰσχρός} ‘ugly, shameful’) → \iwi{AG}{αἰσχύνω} ‘to make ugly, shame’.\footnote{Followed by \citealt{VillanuevaSvensson2024}. Earlier works considered the -\textit{n}- in the suffix -ύνω\iw{AG}{-ῡ́νω} to be simply an extension without further explanation (see \citealt[264]{Palmer1980}).} However, Ancient Greek must have innovated by replacing the process of infixation by \isi{suffixation} followed by the suffix \iwi{PIE}{*-\textit{i̯e/o}-} \is{*-\textit{i̯e/o}-present} in the present stem of the \is{denominative!verb} denominative verbs.\footnote{\citealt[276--277]{Chantraine1947}, \citealt[25]{Lamberterie1990} (“-ῡνω\iw{AG}{-ῡ́νω} < *-ῠν-i̯ω, aor.\ \iwi{AG}{-ῡνα} < *‑ῠν-σα”). On the denominal\is{denominal!verb} nasal suffix \is{nasal present} in the northern Indo-European\is{Northern Indo-European} languages, see \citealt{VillanuevaSvensson2016}.} For a new explanation of the development of this derivational process in Greek and its relationship to the denominative suffix\is{denominative!suffix} \iwi{AG}{-αίνω}, see \citealt{VillanuevaSvensson2024}.

\section{Deadjectival Verbs in Anatolian}
\subsection{Formations in -\textit{nu}-}
\label{sec:nu}

As is well known, the Hittite language has a significant number of \textit{nu}-formations built from adjectives in \iwi{Hitt}{-\textit{u}-}, e.g., \iwi{Hitt}{\textit{parku}-} ‘high’ (< \iwi{PIE}{*\textit{bʰerg̑ʰ-u}-}) → \iwi{Hitt}{\textit{park(a)nu}-} ‘to make high’, cf.\ the root verb \iwi{Hitt}{\textit{parktaru}} (3sg.imp.) ‘must become high’; \iwi{Hitt}{\textit{āššu}-} ‘good’ (< \mbox{ \iwi{PIE}{*\textit{h\textsubscript{1}e(h\textsubscript{1})s-u}-})} → \textit{aš(ša)nu-\textsuperscript{mi}}\iw{Hitt}{\textit{aš(ša)nu}-} ‘to take care of’, cf.\ the \isi{stative}, mediopassive root verb \iwi{Hitt}{\textit{āšš}-} ‘to be dear to’; \iwi{Hitt}{\textit{ḫatku}-} ‘narrow’ (< \iwi{PIE}{*\textit{h\textsubscript{2}edʰgʰ-u}-}) → \textit{ḫatg(a)nu-\textsuperscript{mi}}\iw{Hitt}{\textit{ḫatg(a)nu}-} ‘to make narrow’,  cf.\ the root verb \textit{ḫatk-\textsuperscript{ḫi}}\iw{Hitt}{\textit{ḫatk}-} ‘to close’. However, despite what scholars may claim, it is not always clear whether the \textit{nu}-formation is derived from the \textit{u}-adjective or from the root verb.\footnote{\citet[156]{Luraghi1992} sees \textit{aš(ša)nu-\textsuperscript{mi}}\iw{Hitt}{\textit{aš(ša)nu}-} as deverbal\is{deverbal/deverbative!verb} instead. The etymological rejection of \mbox{\textit{aš(ša)nu-\textsuperscript{mi}}}\iw{Hitt}{\textit{aš(ša)nu}-} made by \citet[216--218]{Kloekhorst2008} is rebutted by \citet[473, fn.\ 10]{Sasseville2021}.} Although some formations may be more transparent than others, it is always possible to argue for the one or for the other direction, which is relevant for the topic of \isi{reanalysis}. Considering solely the deadjectival type\is{deadjectival!verb} from a synchronic point of view, one may speak of an \textit{n}-infix or of a \textit{nu}-suffix with deletion of the \textit{u}-suffix. The latter process is precisely the one which applies to the adjectival \textit{i}-stems, e.g., Hitt.\ \textit{šalli-/-ay-}\iw{Hitt}{\textit{šalli}-} ‘big, great’ (< \iwi{PIE}{*\textit{solh\textsubscript{2/3}-i}-}) → \textit{šall(a)nu-\textsuperscript{mi}}\iw{Hitt}{\textit{šall(a)nu}-} ‘to make big’, \textit{ḫarki-/-ay-}\iw{Hitt}{\textit{ḫarki}-} ‘white’ (< \iwi{PIE}{*\textit{h\textsubscript{2}erg̑-i}-}) → \textit{ḫark(a)nu\nobreakdash-\textsuperscript{mi}}\iw{Hitt}{\textit{ḫark(a)nu}-} ‘to make white’, \textit{parkui-/-ay-}\iw{Hitt}{\textit{parkui}-} ‘pure’ (< \iwi{PIE}{*\textit{perk\textsuperscript{u̯}-i}-})\footnote{For the root, see \citetalias{eDiAna} \href{http://www.ediana.gwi.uni-muenchen.de/dictionary.php?lemma=565}{ID 565}.} → \textit{parkunu-\textsuperscript{mi}}\iw{Hitt}{\textit{parkunu}-} ‘to make pure’ and \textit{mekki-/-ay-}\iw{Hitt}{\textit{mekki}-} ‘many, much’ (< \iwi{PIE}{*\textit{meg̑h\textsubscript{2}-i}-}) → \textit{maknu-\textsuperscript{mi}}\iw{Hitt}{\textit{maknu}-} (\textit{meknu-\textsuperscript{mi}}\iw{Hitt}{\textit{meknu}-}) ‘to make numerous’.\footnote{For a list of examples of the denominal type,\is{denominal!verb} see \citealt[155--156]{Luraghi1992} and \citealt[70--72]{Vanseveren2014}.}

It is appropriate to recall that there are no clear examples of Hittite \textit{a}-stems, which reflect the merger of \iwi{PIE}{*-\textit{o}-} (unless it was syncopated, in which case the stem became consonantal) and \iwi{PIE}{*\textit{-eh\textsubscript{2}}-}, that serve as derivational bases of \textit{nu}-formations. There are a few examples that have been adduced in the literature and the first one is \textit{marš(a)nu-\textsuperscript{mi}}\iw{Hitt}{\textit{marš(a)nu}-} ‘to desecrate’ next to \iwi{Hitt}{\textit{marša}-} ‘false(hood)’.\footnote{This example is left out in \citealt[70--72]{Vanseveren2014}.} For this reason, \citet[29]{Koch1973} finds its interpretation as a derivative from the \textit{a}-stem \iwi{Hitt}{\textit{marša}-} suspect and therefore does not exclude the possibility that it may have been derived from an unattested root verb or even from a prehistoric \textit{u}-stem. Likewise, \citet[239, fn.\ 1]{Oettinger1979} and \citet[157]{Luraghi1992} distance themselves from the nominal base \iwi{Hitt}{\textit{marša}-} and presuppose instead a deverbal\is{deverbal/deverbative!verb} formation from \textit{maršiye/a-\textsuperscript{mi}}\iw{Hitt}{\textit{maršiye/a}-} ‘to be false’ with deletion of the verbal suffix \textit{-iye/a-\textsuperscript{mi}}\iw{Hitt}{-\textit{iye/a}-} (< \iwi{PIE}{*-\textit{i̯e/o}-}).\footnote{Note that there is now even less evidence for a synchronic Hittite class of \isi{stative} verbs in -\textit{e-\textsuperscript{mi}}\iw{Hitt}{-\textit{e}-}, since it has finally been recognized that the cuneiform sign E may have the phonetic value of the glide \iwi{Hitt}{-\textit{i̯}-}; see \citealt[134--161]{Kloekhorst2014} and the subtle warnings in the older literature such as in \citealt[10]{Neu1983} and \citealt[35, 141--142]{Melchert1994}. \citet[225--226]{Oettinger1992} has already pointed out that only \iwi{Hitt}{\textit{ar-ša-ne-e-ši}} and \iwi{Hitt}{\textit{mar-še-e-er}} remain as probative examples. Nonetheless, the unmistakable formation \textit{-iye/a-\textsuperscript{mi}}\iw{Hitt}{\textit{-iye/a}-} in the amply attested \textit{aršaniye/a-\textsuperscript{mi}}\iw{Hitt}{\textit{aršaniye/a}-} as well as the fact that the hapax \iwi{Hitt}{\textit{mar-še-e-er}} could still reflect a stem \textit{maršiye/a-\textsuperscript{mi}}\iw{Hitt}{\textit{maršiye/a}-} despite not being written \iwi{Hitt}{*\textit{mar-ši-e-er}} render the existence of this class synchronically highly questionable.} 

\largerpage
A third possibility that has been overlooked is a deadjectival derivation\is{derivation!deadjectival} \iwi{Hitt}{\textit{maršai(ya)}-} from an \textit{i}-stem \textit{marši-/-ay-}\iw{Hitt}{\textit{marši}-} ‘pertaining to desecration, false’, if it truly is a Hittite lexeme and not a Luwianism.\footnote{Cf.\ \citetalias{eDiAna}  \href{http://www.ediana.gwi.uni-muenchen.de/dictionary.php?lemma=2869}{ID 2869} with further literature.} 

Another possible example for an \textit{a}-stem as derivational base is presented in \citetalias[ Ḫ 87--88]{HW}, where it is argued that \textit{ḫalluwanu-\textsuperscript{mi}}\iw{Hitt}{\textit{ḫalluwanu}-} of uncertain meaning is not derived from \textit{ḫalluwā(i)-\textsuperscript{mi}}\iw{Hitt}{\textit{ḫalluwā(i)}-} ‘to fight’, i.e., ‘to cause to fight’, but rather from \mbox{\textit{ḫalluwa}-}\iw{Hitt}{\textit{ḫalluwa-/ḫalluwaw}-} ‘deep’, i.e., ‘to deepen’ against the previous literature. The argument rests on the understanding of two contexts, which cannot be reviewed within the frame of this article. Moreover, it is worth noting that the adjective \mbox{\textit{ḫalluwa}-}\iw{Hitt}{\textit{ḫalluwa-/ḫalluwaw}-} has a \textit{u}-stem allomorph \textit{ḫalluwaw}- visible in the acc.\ pl.\ \mbox{\textit{ḫalluwawuš}} and \mbox{\textit{ḫalluwāmuš}} (see \citetalias[ Ḫ 84--85]{HW}), which, if this lexeme is indeed the derivational base, may have paved the way for the \textit{nu}-\isi{factitive}, i.e.,\textit{ḫalluwaw}- → \textit{ḫalluwa-n-w}-.\footnote{The hapax \iwi{Hitt}{\textit{maninkuwanut}} (CHDS 3.3, 7') is not a relevant example because its derivational base is not (properly) attested (see \citetalias[ L--N 170]{CHD} \iwi{Hitt}{*\textit{maninkuwa}-} ‘near’). Therefore, the stem formation of the base is unknown; cf.\ \citealt[246]{Oettinger1979} (“\iwi{Hitt}{\textit{maninkui}-}”).} A similar problematic case concerns Hitt.\ \textit{ḫatuganu-\textsuperscript{mi}}\iw{Hitt}{\textit{ḫatuganu}-} ‘to frighten’, which can be derived from the adjective \textit{ḫatuga}-\iw{Hitt}{\textit{ḫatuga}-/\textit{ḫatugi}-} (or \textit{ḫatugi}-) ‘frightful’ \citep[155]{Luraghi1992} or from the root verb \textit{ḫatuk-\textsuperscript{mi}}\iw{Hitt}{\textit{ḫatuk}-} ‘to spread fear (intr.)’ (\citealt{Koch1973}: 19f.; \citetalias[ Ḫ 529]{HW}), both semantically adequate, i.e., ‘to cause to spread fear’ or ‘to make frightful’.\footnote{On the few spellings \iwi{Hitt}{\textit{ḫatuganu}-} for expected \iwi{Hitt}{*\textit{ḫatuknu}-}, cf.\ the similar situation of \iwi{Hitt}{\textit{zalukanu}-} attested beside the variant \iwi{Hitt}{\textit{zaluknu}-} (\citealt[1027--1028]{Kloekhorst2008}).} Moreover, the adjective \textit{ḫatuga}-\iw{Hitt}{\textit{ḫatuga}-/\textit{ḫatugi}-} shows an \textit{i}-stem allomorph \textit{ḫatugay}- in, e.g., abl.\ \mbox{\textit{ḫatugayaz},} nom.\ pl.\ \textit{ḫatugayeš} (\textit{ḫa-tu-ga-e-eš}) and acc.\ pl.\ \textit{ḫatugauš} (\citetalias[ Ḫ 525--529]{HW}), which may have been relevant for the derivational process, if the adjective turns out to be the proper base. Be that as it may, it is safe to posit that there is no unambiguous  synchronic evidence that would show that the Hittite \isi{factitive} \textit{nu}-suffix can derive verbs from \textit{a}-stems.\footnote{As pointed out by \citet[483]{Sasseville2021}, since the Common Anatolian thematic suffix \iwi{PAn}{*-\textit{o}-} was syncopated in post-tonic position and the formation transferred to another stem class, e.g., \textit{u}-stems, it is always possible that a formation in \iwi{PAn}{*-\textit{nu}-} was originally created at the time when the thematic suffix was still present. Unfortunately, there is as of yet no positive evidence for this claim.}

In the smaller Luwian corpus, we find the several cases of the denominal type\is{denominal!verb} in /-nuwa-\textsuperscript{i}/\iw{CLuw}{-\textit{nuwa}-}\iw{HLuw}{-\textit{nuwa}-} (cf.\ \citealt[471--477, 484--485]{Sasseville2021}) derived from \textit{u}-adjectives, e.g.,  (\ref{ex:Luw}).

\begin{exe}
\ex \label{ex:Luw} \begin{xlist}
    \ex CLuw.\ \iwi{CLuw}{\textit{wāšu}-},  HLuw.\ \iwi{HLuw}{\textit{wa/i-su}-} ‘good’ (< *\textit{h\textsubscript{1}u̯os-u}-)\iw{PIE}{*\textit{h\textsubscript{1}u̯ós-u-/h\textsubscript{1}u̯és-u}-} → HLuw.\ \iwi{HLuw}{\textit{u-sa-nú}-} /usnuwa-\textsuperscript{i}/ ‘to make good, bless’ 
    \ex Luw.\ */assu-/\iw{CLuw}{*\textit{assu}-}\iw{HLuw}{*\textit{assu}-} ‘good’ (= Hitt.\ \iwi{Hitt}{\textit{āššu}-} < \iwi{PIE}{*\textit{h\textsubscript{1}e(h\textsubscript{1})s-u}-}) → HLuw.\ \iwi{HLuw}{\textit{a-su-nu}-} /ass(u)nuwa-\textsuperscript{i}/ ‘to make good (?)’ (cf.\ Lyd. \textit{f-ašνo-\textsuperscript{d}}\iw{Lydian}{\textit{fašνo}-} ‘to make good (?)’)\footnote{See \citealt[479--480]{Sasseville2021}.}
    \ex HLuw.\ */habanzu-/ ‘loyal’ (< \iwi{PIE}{*\textit{h\textsubscript{2}op-ent-i̯u}-}\footnote{Cf.\ \citetalias{eDiAna}  \href{http://www.ediana.gwi.uni-muenchen.de/dictionary.php?lemma=1580}{ID 1580} for various ways of interpreting the morphological chain.}) → HLuw.\ \iwi{HLuw}{\textit{ha-pa-za-nu-wa}-} /habanznuwa-\textsuperscript{i}/ ‘to make loyal’ (cf.\  CLuw.\ \includegraphics[]{Glossenkeil}\iwi{CLuw}{\textit{ḫapanzuwalatar/n}-} ‘loyalty’ and HLuw.\ \iwi{HLuw}{\textit{ha-pa-zú+ra/i-wa/i}-} /habanzulawa(/i)-/ ‘devoted, devotion (?)’)
    \end{xlist}
    \end{exe}
    
    From a synchronic perspective, these formations can be interpreted either as \textit{n}-infixed or as \textit{nu}-suffixed\is{suffixation} via deletion of the adjectival suffix -\textit{u}-, although the former is expected to be the inherited process. The rest of the deadjectival derivations\is{derivation!deadjectival}  with the Luwian suffix /-nuwa-\textsuperscript{i}/\iw{CLuw}{-\textit{nuwa}-} \iw{HLuw}{-\textit{nuwa}-} are built on mutated\is{\textit{i}-mutation} stems that can go back to either a thematic \textit{o}-stem or an \textit{i}-stem, e.g., CLuw.\ \iwi{CLuw}{\textit{ura/i}-}, HLuw.\ \iwi{HLuw}{MAGNUS\textit{+ra/i}-} ‘big, great’ (< \iwi{PIE}{*\textit{h\textsubscript{1}ur-ó}-})\footnote{See \citetalias{eDiAna}  \href{http://www.ediana.gwi.uni-muenchen.de/dictionary.php?lemma=3247}{ID 3247}.} → CLuw.\ \textit{uranuwa-\textsuperscript{i}}\iw{CLuw}{\textit{uranuwa}-}, HLuw.\ \iwi{HLuw}{MAGNUS\textit{+ra/i-nú-wa/i}-} /uranuwa-\textsuperscript{i}/ ‘to make great’; CLuw.\ \iwi{CLuw}{\textit{ḫalāl(i)}-} ‘pure’ (Semitic borrowing) → \textit{ḫalalanuwa-\textsuperscript{i}}\iw{CLuw}{\textit{ḫalalanuwa}-} ‘to make pure’; CLuw.\ \iwi{CLuw}{*\textit{zantal(i)}-} ‘low’ (< \mbox{\iwi{PIE}{*\textit{\kinvbreve m̥t-ó-lo}-})} → CLuw.\   \includegraphics[]{Glossenkeil}\textit{zantalanuwa-\textsuperscript{i}}\iw{CLuw}{\textit{zantalanuwa}-} ‘to make low, diminish’.\footnote{The lack of a clear example going back to a PIE \textit{i}-stem is simply due to chance.} As can be gathered from these examples, there is no process of deletion here. Instead, the suffix \textit{\nobreakdash-nuwa\nobreakdash-\textsuperscript{i}}\iw{CLuw}{-\textit{nuwa}-} \iw{HLuw}{-\textit{nuwa}-} is added to the thematic vowel \iwi{CLuw}{-\textit{a}-}\iw{HLuw}{-\textit{a}-} (< \iwi{PIE}{*-\textit{o}-}), which is the regular derivational vowel in the Luwian \textit{i}-mutated\is{\textit{i}-mutation} paradigm. In comparison with Hittite, however, one may posit that the \isi{factitive} suffix was able to reach the thematic stems in Luwian only because of their merger with the \textit{i}-stems and that the process of deleting the suffix -\textit{i}-, assuming that it existed in Pre-Luwian, was given up, i.e., *-\textit{i}- (with deletion) + *-\textit{nu}- → *-\textit{o-/-i}- + \textit{\nobreakdash-nu-} = *-\textit{o-nu}-. On comparative grounds, adjectival \textit{u}-stems inherited the \textit{n}-infixation and, as far as we can tell, preserved it.

The next question concerns the path of the extension from \textit{u}-adjectives as the derivational base to \textit{i}-adjectives and the reason behind the \isi{reanalysis} of the infixation process as \isi{suffixation}. The two following scenarios require  infixation to not only be genuine but also older: The first explanation is that old \textit{u}-adjectives were gradually replaced  by younger \textit{i}-adjectives, which motivated the reinterpretation of the \isi{factitive} formation from infixation to \isi{suffixation} via deletion of the suffix, i.e., CVC-\textit{u}- : CVC-\textit{n-u}- reanalyzed\is{reanalysis} as CVC-\textit{i}- (ousting the \textit{u}-stem) : CVC-\textit{nu}- (\citealt[35]{Koch1973}, \citealt[27]{Tucker1981}). On the other hand, it has been claimed that the \isi{reanalysis} of the derivational process simply happened due to the decline of the infixation process altogether, which led to a new suffix *-\textit{nu}- creating \isi{factitive}s via deletion of the adjectival suffix *-\textit{u}- and then *-\textit{i}- (cf.\ \citealt[245--247]{Oettinger1979}, \citeyear[157]{Oettinger2022}).\footnote{However, the claim found in, e.g., \citealt[246--247]{Oettinger1979} that Hittite has stems in \iwi{Hitt}{-\textit{ui}-} (< \iwi{PIE}{*-\textit{u̯-ih\textsubscript{2}}}) is considered outdated here; see \citealt[829]{Kloekhorst2008}.} The reason why the \textit{i}-stems would have been targeted would simply be due to their similarity to the \textit{u}-stems. In the end, the second scenario has the advantage of better explaining why it is the \textit{i}-stems that followed the path of the \textit{u}-stems without having to recur to the loss of an original \textit{u}-adjective.

\subsection{New ways of forming verbs from Anatolian \textit{u}-adjectives}
\label{sec:uadj}

The suffix \iwi{Hitt}{-\textit{nu}-} is not the only one that can create deadjectival verbs.\is{deadjectival!verb} At some point, Hittite started using the \isi{factitive} suffix -\textit{aḫḫ-\textsuperscript{i}}\iw{Hitt}{-\textit{aḫḫ}-}, which originates in stems in \iwi{PIE}{*\textit{-eh\textsubscript{2}}-}, to derive \textit{u}-adjectives via their oblique stem, e.g., \textit{tēpu-/-aw-}\iw{Hitt}{\textit{tēpu}-} ‘small’ → \textit{tepawaḫḫ-\textsuperscript{i}}\iw{Hitt}{\textit{tepawaḫḫ}-} ‘to make small, humble’ (vs.\ \textit{tepnu-\textsuperscript{mi}}\iw{Hitt}{\textit{tepnu}-}), \textit{idālu-/-aw-}\iw{Hitt}{\textit{idālu}-} ‘bad, evil’ → \textit{idalāwaḫḫ-\textsuperscript{i}}\iw{Hitt}{\textit{idalāwaḫḫ}-} ‘to treat badly (tr.), to act badly (intr.)’.\footnote{There is a widespread view that \iwi{Hitt}{-\textit{aḫḫ}-} was mainly derived from stems in \iwi{Hitt}{-\textit{ant}-} via deletion of the suffix (see \citealt[240--241]{Oettinger1979}, \citealt[147]{Vanseveren2014}), which would be possible as a secondary phenomenon, but is unconvincing  since all examples can instead be derived from the base in \iwi{Hitt}{-\textit{a}-}, which was often extended with the highly productive\is{productivity} denominal\is{denominal!verb} suffix \iwi{Hitt}{-\textit{ant}-}. See \citealt[284--285]{Sasseville2015} for this criticism.} There is also one potential example in Palaic, i.e., \iwi{Palaic}{\textit{wāšu}-} ‘good’ → \iwi{Palaic}{\textit{wašūḫ}} (2sg.imp.) ‘to act well, be good.\footnote{See \citealt[536, fn.\ 8]{Sasseville2021}.} Nonetheless, the number of examples is very low, so that the rivalry between \iwi{Hitt}{-\textit{aḫḫ}-} and \iwi{Hitt}{-\textit{nu}-} regarding adjectival \textit{u}-stems is subtle. \citet[282]{Shatskov2020} claims that the two Hittite examples mentioned above (though re-addressing only \textit{tepawaḫḫ-\textsuperscript{i}}\iw{Hitt}{\textit{tepawaḫḫ}-}) were created by non-native speakers who did not grasp  Hittite morphology, but this is not persuasive. Although there is a strong interference from Luwian native speakers in New Hittite, this is certainly not the case here because the Luwian \isi{factitive} suffix was restricted to \isi{conversion} and, therefore, there was no such model for deriving \textit{u}-stems in their own language.\footnote{On Luwian interference in Hittite, see \citealt{Melchert2005}, \citealt{Rieken2006}, \citealt{Hout2006}, \citealt[15--73]{Yakubovich2010}.} The Hittite \isi{factitive} suffix \iwi{Hitt}{-\textit{aḫḫ}-} simply became highly productive\is{productivity} and extended its derivational scope to all kinds of nominal stems.

\citet[29--30]{Koch1973} notes that the extension of \iwi{Hitt}{-\textit{aḫḫ}-} to the adjectival \textit{i}-stems is largely restricted to synchronically non-ablauting \textit{i}-stems, e.g., \iwi{Hitt}{\textit{nakkī}-}\footnote{For the exceptional origin of this particular \textit{i}-stem, see \citealt{Widmer2005}.} ‘heavy, important’ → \textit{nakkiyaḫḫ-\textsuperscript{i}}\iw{Hitt}{\textit{nakkiyaḫḫ}-} ‘(intr.) to become a concern to’, \iwi{Hitt}{\textit{šannapili}-} ‘empty’ (but cf.\ acc.\ pl.\ [\textit{ša-an-n}]\textit{a-pí-la-a-uš}; \citetalias[ Š 159]{CHD}) → \textit{šannapilaḫḫ-\textsuperscript{i}}\iw{Hitt}{\textit{šannapilaḫḫ}-} ‘to empty’, \iwi{Hitt}{\textit{šanezzi(ya)}-} ‘pleasant’ → \textit{šanezziyaḫḫ-\textsuperscript{i}}\iw{Hitt}{\textit{šanezziyaḫḫ}-} ‘to indulge oneself’, \iwi{Hitt}{\textit{šarazzi(ya)}-} ‘upper’ → \textit{šarazziyaḫḫ-\textsuperscript{i}}\iw{Hitt}{\textit{šarazziyaḫḫ}-} ‘to make prevail’. Regarding the suffix \iwi{Hitt}{-\textit{zzi(ya)}-}, we are dealing with an originally thematic suffix *´-\textit{ti̯o}-\iw{PIE}{*-\textit{ti̯o}-} that secondarily underwent contraction or syncope of the thematic vowel, thus inflecting sometimes like an \textit{i}-stem. A similar explanation applies to stems in \iwi{Hitt}{-\textit{ili}}, whose suffix originates in \iwi{PIE}{*-\textit{í-lo}-} with syncope of the thematic vowel in post-tonic position followed by a secondary transfer of the stem class.\footnote{See \citealt{Rieken2008}.} These data indicate that the Hittite \isi{factitive} suffix \mbox{\iwi{Hitt}{-\textit{aḫḫ}-}} was still closely associated with the \textit{a}-stems (\iwi{PIE}{*-\textit{eh\textsubscript{2}}-}, \iwi{PIE}{*-\textit{o}-}) and, after syncope of *´-\textit{o}-, with the consonantal stems, after which it was even extended to the (secondary) non-ablauting \textit{i}-stems.\is{derivation!stem-based} An exception to this pattern is \textit{šuppi-/-ay-}\iw{Hitt}{\textit{šuppi}-} ‘pure’ → \textit{šuppiyaḫḫ-\textsuperscript{i}}\iw{Hitt}{\textit{šuppiyaḫḫ}-} ‘to purify’, which would show that the \isi{factitive} suffix in the end did reach the genuine ablauting \textit{i}-stems, conflicting once more with the \isi{factitive} suffix \iwi{Hitt}{-\textit{nu}-}.

Overall, the Hittite \isi{factitive} suffix \iwi{Hitt}{-\textit{aḫḫ}-} was so productive\is{productivity} that it could be derived from any substantival or adjectival nominal stem, whereas the \isi{factitive} suffix \iwi{Hitt}{-\textit{nu}-}, based on the attested data, was confined to adjectival \textit{u}- and \textit{i}-stems.\footnote{It is not entirely clear whether Hittite \iwi{Hitt}{-\textit{nu}-} could derive verbs from substantival bases; cf.\ \citealt[71]{Vanseveren2014}, who lists three examples: \textit{ešḫarnu-\textsuperscript{mi}}\iw{Hitt}{\textit{ešḫarnu}-} ‘to make bloody’, \textit{ēdriyanu-\textsuperscript{mi}}\iw{Hitt}{\textit{ēdriyanu}-} ‘to feed’ and \textit{šaminu-\textsuperscript{mi}}\iw{Hitt}{\textit{šaminu}-} ‘to burn (trans.)’, although a deverbal derivation\is{deverbal/deverbative!verb} based on the denominative verb\is{denominative!verb}  \textit{ēdriye/a-\textsuperscript{mi}}\iw{Hitt}{\textit{ēdriye/a}-} ‘to feed’ is another viable option (thus \citealt[175, fn.\ 10]{HoffnerMelchert2008}; cf.\ \citealt[155]{Luraghi1992}) which deserves further investigation. The same ambiguity is found in \iwi{Hitt}{KAL-\textit{tarnuškanzi}} /hadugadarnuskanzi/ (KBo 12.109, 7'), which may be derived from the abstract noun in \textit{-ātar} (i.e., \iwi{Hitt}{\textit{ḫatugātar}}) or from an unattested denominative\is{denominative!verb} in *\iwi{Hitt}{-\textit{iye/a}-} with deletion of the suffix. The derivational base of \textit{šaminu-\textsuperscript{mi}}\iw{Hitt}{\textit{šaminu}-} is problematic and thus debatable (cf.\ \citetalias[ Š 118]{CHD}, \citealt{Luraghi1992}: 158). One last possible example of a denominal\is{denominal!verb} case is the hapax \textit{aimpanu\nobreakdash-\textsuperscript{mi}}\iw{Hitt}{\textit{aimpanu}-} ‘to burden’ (\citetalias[ A 48]{HW}), if it is derived from the noun \textit{aimpa-/impa-} c.\ \iw{Hitt}{\textit{aimpa-/impa-} c.} ‘burden’ instead of from the denominative verb\is{denominative!verb} *\textit{(a)impā(i)-\textsuperscript{mi}}\iw{Hitt}{*\textit{(a)impā(i)}-} ‘to burden’ (attested solely as a verbal noun \iwi{Hitt}{\textit{impawar}} and as an \isi{imperfective} \textit{impaiške/a-\textsuperscript{mi}}\iw{Hitt}{\textit{impaiške/a}-} according to \citetalias[ I 56]{HW}). Despite the occasional ambiguity, it does appear as if a desubstantival derivation for \iwi{Hitt}{-\textit{nu}-} was slowly becoming possible, perhaps through reinterpretation of the derivational base.} The reason behind this distribution may be the nature of the process itself: Whereas the former is simply suffixed\is{suffixation} to the nominal stem in Hittite, the latter uses either infixation or \isi{suffixation} with deletion of the nominal suffix of the base, which must have been morphologically less stable.

There is a view that verbs in -\textit{uwe/a-\textsuperscript{mi}}\iw{Hitt}{-\textit{uwe/a}-}, which are mostly deradical\is{derivation!deradical} to roots ending in \iwi{Hitt}{-\textit{u}-} derived with the suffix \textit{-iye/a-\textsuperscript{mi}}\iw{Hitt}{-\textit{iye/a}-}, may also be denominal\is{denominal!verb} \is{deadjectival!verb} to \textit{u}-adjectives, e.g. \textit{ḫuišu-/-aw-}\iw{Hitt}{\textit{ḫuišu}-} ‘alive, raw’ → \textit{ḫuišuwe/a-\textsuperscript{mi}}\iw{Hitt}{\textit{ḫuišuwe/a}-} ‘to live, stay alive’. Dismissing this morphological interpretation, as well as that of a \isi{stative} in -\textit{e\nobreakdash-\textsuperscript{mi}}\iw{Hitt}{-\textit{e}-} because of the lack of parallels, \citet[69]{Jasanoff2023} analyzes this example as a thematized\is{thematization} \isi{*-\textit{u}-present},\iw{PIE}{*-\textit{u}-} comparing the reconstructed PIE *\textit{g\textsuperscript{u̯}íh\textsubscript{3}u̯e/o}-\iw{PIE}{*\textit{g\textsuperscript{u̯}ih\textsubscript{3}-u̯é/ó}-} ‘live’.\footnote{See also \citealt[273--274\ with fn.\ 33]{VillanuevaSvensson2021}.} As per Jasanoff, the evidence for a denominal verb\is{denominal!verb} in \textit{-iye/a-\textsuperscript{mi}}\iw{Hitt}{-\textit{iye/a}-} derived from a \textit{u}-adjective, albeit formally possible, is lacking independent support. However, what has not been considered is the formal possibility of a deradical verb\is{derivation!deradical} in \textit{-iye/a-\textsuperscript{mi}}\iw{Hitt}{-\textit{iye/a}-} derived from a \textit{u}-extended root. This topic as well as this example is addressed again in \sectref{sec:u-ext} of this paper.

Turning to the Luwian language, \citet[190--199]{Sasseville2021} suggests the existence of a denominal \is{denominal!verb} suffix /-(a)u-\textsuperscript{di}/\iw{CLuw}{/-(a)u-/}\iw{HLuw}{/-(a)u-/} used to derive \textit{u}-stems in order to explain the strange morpheme \iwi{CLuw}{/-(a)u-/}\iw{HLuw}{/-(a)u-/} found in some Luwian verbs, although a coherent picture of the data could not be successfully provided. \citet{MelchertYakubovich2022} take a different approach and convincingly show that the Luwian endings \iwi{CLuw}{-\textit{unta}}\iw{HLuw}{-\textit{unta}} and \iwi{CLuw}{-\textit{aunta}}\iw{HLuw}{-\textit{aunta}} do not stand for the 3pl.\ preterit, as previously thought, but for the 1pl.\ preterit ending. Their findings therefore eliminate the evidence for a special denominal suffix \is{denominal!verb} /-(a)u-\textsuperscript{di}/\iw{CLuw}{/-(a)u-/} \iw{HLuw}{/-(a)u-/} in Luwian. In other words, there is currently no other way to derive adjectival \textit{u}-stems in Luwian except with the suffix \mbox{/-nuwa-\textsuperscript{i}/}\iw{CLuw}{-\textit{nuwa}-}.\iw{HLuw}{-\textit{nuwa}-} As already mentioned, the \isi{factitive} suffix \iwi{PIE}{*-\textit{eh\textsubscript{2}}-} is not a potential candidate either because it is restricted to \isi{conversion} in Luwian. On the other hand, the Luwian suffix /-nuwa-\textsuperscript{i}/\iw{CLuw}{-\textit{nuwa}-}\iw{HLuw}{-\textit{nuwa}-} could derive verbs from \textit{a/i-}stems (< \iwi{PIE}{*-\textit{o}-}, \iwi{PIE}{*-\textit{i}-}) and perhaps even from \textit{a}-stems (< *-\textit{eh\textsubscript{2}}-), thus being more susceptible to encroach on the \isi{factitive} suffix \iwi{PIE}{*-\textit{eh\textsubscript{2}}-}. In the end, this yields a situation that is the opposite of the Hittite one.\footnote{Cf.\ \citealt[471--477]{Sasseville2021}. However, due to the size of the Luwian corpus and the lack of data, it is not always clear whether the derivational base was mutated\is{\textit{i}-mutation} (i.e., \iwi{PIE}{*-\textit{o}-}, \iwi{PIE}{*-\textit{i}-}) or not (\iwi{PIE}{*-\textit{eh\textsubscript{2}}-}).}

Turning to the Anatolian languages of the 1\textsuperscript{st} millennium BC, Lycian has basically eliminated nominal \textit{u}-stems, transferring them to the paradigm of the \textit{a}-stems (< \iwi{PIE}{*-\textit{eh\textsubscript{2}}-}), so that it does not offer any evidence whatsoever for the current topic. Lydian, however, seems to have had highly productive\is{productivity} ``\textit{u}''-stems, so that we might find denominal verbs\is{denominal!verb} \is{deadjectival!verb} derived from adjectival \textit{o}-stems (< PIE \iwi{PIE}{*-\textit{u}-}), e.g., \iwi{Lycian}{\textit{taso}-} ‘strong, heavy (\textit{vel sim}.)’ (= Hitt.\ \iwi{Hitt}{\textit{daššu}-} ‘heavy’, CLuw.\ \iwi{CLuw}{\textit{kwanzu}-} ‘id.’) → \textit{taso-\textsuperscript{d}} \iw{Lycian}{\textit{taso}-} ‘to impose’.\footnote{\citealt[180--190]{Sasseville2021}; differently \citealt[174, fn.\ 20]{GarnierSagot2020}.} Taking this example seriously, Lydian appears to derive denominal verbs\is{denominal!verb} in -\textit{o-\textsuperscript{d}}\iw{Lycian}{-\textit{o}-} of the lenited \isi{\textit{mi}-conjugation} from adjectival stems in -\textit{o}- via \isi{conversion}, although the mechanism in this case is more likely to be secondary, rather than inherited.\footnote{See \citealt[203--205]{Sasseville2021}.}

\section{Reanalysis and \textit{u}-extensions}
\label{sec:u-ext}

According to the analysis of \citet[60--108]{Strunk1967}, there is a difference between \textit{n}-infixed verbs to roots ending in -\textit{u} (e.g., \iwi{PIE}{*\textit{\kinvbreve leu̯}-} → *\textit{\kinvbreve l̥-né-u̯-/-n-u}-´)\iw{PIE}{*\textit{\kinvbreve l̥-né-u̯-/\kinvbreve l̥-nu}-´} and \textit{n}-infixed verbs to roots ending in a laryngeal, in which the -\textit{u}-, unlike the laryngeal, is not part of the root; compare Ved.\ \iwi{Sanskrit}{\textit{śr̥ṇā́ti}} ‘to break apart’, \iwi{Sanskrit}{\textit{śīrṇá}-}, aor.\ \iwi{Sanskrit}{\textit{aśarīt}}, \iwi{Sanskrit}{\textit{śīrti}-} ‘breaking’ with Ved.\ \iwi{Sanskrit}{\textit{vr̥ṇóti}} ‘to enclose’‚ \iwi{Sanskrit}{\textit{vr̥tá}-}, aor.\ \iwi{Sanskrit}{\textit{ā́var}}, \iwi{Sanskrit}{\textit{vr̥ti}-} ‘fencing’. Although Strunk argues that the origin of “\textit{nu}-presents” \is{nasal present} is to be found in roots that end in -\textit{u},  apart from \iwi{PIE}{*\textit{\kinvbreve leu̯}-} ‘to hear’ and \iwi{PIE}{*\textit{ki̯eu̯}-} (Gk.\  \iwi{AG}{κίνυται} ‘to be in motion’, \iwi{AG}{ἔσσυτο}, OInd.\ \iwi{Sanskrit}{\textit{cyuti}-} ‘rapid motion’), he only provides examples of these so-called \textit{u}-extended roots, e.g., \iwi{Sanskrit}{\textit{dabhnóti}} (!), \iwi{Sanskrit}{\textit{vr̥ṇóti}} ‘defends’ (cf.\ Gk.\ \iwi{AG}{ἔρυσθαι}), \iwi{Sanskrit}{\textit{vr̥ṇóti}} ‘to enclose’ (cf.\ Gk.\ \iwi{AG}{εἰλύω} ‘id.’).  This demonstrates at the very least that there is a certain connection between \textit{nu}-verbs and \textit{u}-extensions.

The phenomenon of \textit{u}-extension in the Indo-European languages is an obscure topic, especially vis-à-vis the so-called “\textit{u}-present” \is{*-\textit{u}-present}  suffix.\footnote{Cf.\ \citealt[60--61]{Strunk1967}, \citealt[231--233 w.\ fn.\ 4]{Klingenschmitt1982}, \citealt[211--217]{Steer2015}, \citealt{Jasanoff2023}; cf.\ also \citealt[53]{Watkins1969} on “\textit{u}-endings”. The “\textit{u}-present”\is{*-\textit{u}-present} explanation is used for Hittite \textit{tarḫu-\textsuperscript{mi}}\iw{Hitt}{\textit{tarḫu}-} and \textit{lāḫū-\textsuperscript{ḫi}}\iw{Hitt}{\textit{lāḫū}-} by \citet[222--223., 424]{Oettinger1979}. Contra \citet[67 with fn.\ 25]{Jasanoff2023}, Hitt.\ \textit{tarḫu-\textsuperscript{mi}}\iw{Hitt}{\textit{tarḫu}-} can only be interpreted as a root verb of the \isi{\textit{mi}-conjugation} and the lack of assimilation to \iwi{Hitt}{*\textit{tarruzzi}} is due to the generalization of the weak stem, where the accent was not on the full-grade root, across the paradigm (cf.\ Hitt.\ \textit{walḫ-\textsuperscript{mi}}\iw{Hitt}{\textit{walḫ}-} ‘to strike’ {<}{<} \iwi{PIE}{*\textit{u̯élh\textsubscript{2}-/u̯l̥h\textsubscript{2}}-´}). The two attestations \iwi{Hitt}{\textit{tar-ḫu-id-du}} (KUB 36.75 iv 10') and \iwi{Hitt}{\textit{tar-ḫu-e-zi}} (KBo 38.126, r.col. 10'), which he uses as evidence for an inherited thematic root verb,\is{thematic!present} are not decisive: Both attestations recall the occasional marking of an extra \textit{-i-}/-\textit{e-} vowel following a -\textit{u}-, in this case combined with a marker of the fortis quality of the ending in the first example breaking the consonant cluster, i.e., \mbox{[tárhhʷ\textsubscript{ə}ttu]} and \mbox{[tárhhʷ\textsubscript{ə}tsi]} respectively; cf.\ \iwi{Hitt}{\textit{tar-aḫ-ḫu-i-iš-ke/a}-} \mbox{[tarhhʷ\textsubscript{ə}ské/á-]} (KBo 61.221 obv.\textsuperscript{?} r.col. 2'), \iwi{Hitt}{\textit{pa-aḫ-ḫu-˹i˺-na-li-az}} (VBoT 58 iv 36 to \iwi{Hitt}{\textsuperscript{DUG}\textit{paḫḫunal(l)i}-} n.\ ‘fire container’), \iwi{Hitt}{\textit{wa-at-ku-it-ta}} (KUB 30.67, 6') and \iwi{Hitt}{\textit{wa-at-ku-at-ta}} (KBo 13.137, 11' to \textit{watku-\textsuperscript{mi}}\iw{Hitt}{\textit{watku}-} ‘to jump’; pace \citealt[337, fn.\ 161]{Oettinger1979}) and similarly with \textit{i}-glide insertion between \textit{u} and another vowel \iwi{Hitt}{[\textit{te}]\textit{\nobreakdash-pa\nobreakdash-u\nobreakdash-i\nobreakdash-e\nobreakdash-eš\nobreakdash-ta}} (KBo 58.8+ rev. 23'), \iwi{Hitt}{\textit{ša-an-ḫu-u-i-en}} ‘we swept’ (KBo 38.257, 19'), [\textit{ar-k}]\textit{u-i-an-zi} (KBo 37.70, 6'; \citetalias[ A 311]{HW}) and [\textit{a}]\textit{r-ku-ia-a}[\textit{n}\textsuperscript{?}\textit{-zi}] (KBo 40.328, 10' to \textit{arku-\textsuperscript{mi}}\iw{Hitt}{\textit{arku}-} ‘to chant’ attested beside the denominative\is{denominative!verb} \textit{arkuwā(i)-\textsuperscript{mi}}\iw{Hitt}{\textit{arkuwā(i)}-} ‘to pray, chant’). On the other hand, if one wishes to take these spellings with \textit{-i-/-e-} seriously, a deradical verb\is{derivation!deradical} in \iwi{PIE}{*-\textit{i̯e/o}-} \is{*-\textit{i̯e/o}-present} derived from the \textit{u}-extended root with the loss of the \textit{i̯}-glide remains a viable alternative. In any case, assuming the existence of an additional verbal stem formation which is not securely attested in the first place based on these two forms is highly dubious.} \citet{GarciaRamon2021} has recently attempted to differentiate the two by using a semantic approach. He posits that the morpheme -\textit{u}- should be seen as an extension, whenever there is no semantic distinction between the root verb with -\textit{u}- and the one without. However, if the two happen to differ semantically, he claims that the morpheme acts as an aspectual marker (\isi{imperfective}/\isi{perfective}), insofar as it appears only in the present stem. If it is not restricted to the present stem, the morpheme is seen as a marker of \textit{Aktionsart}.\footnote{Unfortunately, the Hitt.\ verb †\textit{tarḫ-\textsuperscript{mi}}\iw{Hitt}{†\textit{tarḫ}-} ‘to overcome’ used by \citet[157--158]{GarciaRamon2021} is a ghostword (outdated \citealt[222--223]{Oettinger1979}), since only a \textit{u}-extended root verb is truly attested in Hittite, i.e., \textit{tarḫu-\textsuperscript{mi}}\iw{Hitt}{\textit{tarḫu}-}; see already \citet[837]{Kloekhorst2008} who convincingly shows that all spellings of \textit{tar-Vḫ-} should be read \textit{tar-uḫ-} not **\textit{tar-aḫ-}.} This hypothesis definitely needs more evidence from a larger set of data, which we cannot pursue in its entirety here. The present chapter will primarily be concerned with the morphological aspect of the problem, taking Anatolian “\textit{u}-verbs” as one category. Nevertheless, attention will also be paid to the semantics of the formations in question.

\largerpage
In the Anatolian languages, there are adjectival \textit{u}-stems attested next to root verbs, to formations in *-\textit{nu}- as well as to \textit{u}-extended verbs, all built on the same root, and it is tempting to see a connection.\ A similar picture has already been described for Hittite and Indo-Iranian.\ \citet[235]{Koch1980} argues based on Old Indic \iwi{Sanskrit}{\textit{dabhnóti}} ‘to harm, deceive’ next to \iwi{Sanskrit}{\textit{ádbhuta}-} ‘wonderful’ that the deadjectival \is{deadjectival!verb} \textit{n}-infixed formation could at some point be mistaken by native speakers as an \textit{n}-infixed verb.\is{nasal present} In this way, the morpheme -\textit{u}- was reanalyzed as part of the root, i.e., \iwi{PIE}{*\textit{debʰu}-} or \iwi{PIE}{*\textit{dbʰeu̯}-}, and, in a similar way, the Old Indic 8\textsuperscript{th} class with, e.g., \iwi{Sanskrit}{\textit{kru}-} ‘to do’ may have originated as, e.g.,  \iwi{Sanskrit}{\textit{kr̥}-} ‘to do’ → \textit{kr̥ṇóti/kr̥ṇvánti}\iw{Sanskrit}{\textit{kr̥ṇóti, kr̥ṇvánti}} → \isi{reanalysis} of the root as \iwi{Sanskrit}{\textit{kru}-} ‘to do’ → \textit{karoti/kurvanti}\iw{Sanskrit}{\textit{karoti, kurvanti}} (cf.\ \citealt[265--268]{ThumbHauschild1959} on the problems of this class and its examples).\footnote{Contra \citet{Hoffmann1976karoti}, who proposes a special phonological development to derive the \textit{u}-verb from the \textit{nu}-verb, showing on philological grounds that the former has fully replaced the latter over time. One should also note the ambiguity in another example of the same class, i.e., \iwi{Sanskrit}{\textit{tan}-} ‘to stretch’ → \textit{tanóti/tunvánti}\iw{Sanskrit}{\textit{tanóti, tunvánti}} < *\textit{tn̥-néu̯-ti/tn̥-nu̯-énti} or < \iwi{PIE}{*\textit{ten-éu̯-ti/t(e)n-u̯-énti}}. Assuming the first formation to be original, the morpheme boundary between root and suffix may have been redefined and the derivational process reinterpreted, i.e., as \textit{tan-ó-ti/tun-v-ánti}; cf.\ \citealt[265]{ThumbHauschild1959} with further literature. Differently, \citet[63]{Jasanoff2023} takes the middle “\textit{u}-present” \is{*-\textit{u}-present} to be original (Ved.\ \textit{tanute})\iw{Sanskrit}{\textit{tanóti, tunvánti}} and the “\textit{nu}-present” \is{nasal present} to be a \isi{backformation} thereof.} Independent evidence for the validity of the reanalyzed root \iwi{PIE}{*\textit{dbʰeu̯}-} comes from Old Avestan \iwi{Avestan}{\textit{dǝbaoman}-} ‘deception (\textit{vel sim}.)’ with the full grade anchored before -\textit{u}-. According to Koch, the cause behind this \isi{reanalysis} was the disappearance of the \textit{u}-adjective, i.e., the ultimate derivational base, from the language or languages concerned. However, it is questionable whether this is truly the case and, therefore, it is the goal of the last section of this paper to test Koch’s scenario against a new set of Anatolian data.

The first root examined here is \iwi{PIE}{*\textit{h\textsubscript{1}u̯es}-} ‘good’: A \isi{stative} verb can be posited on the basis of the CLuw.\ (though attested in Hittite context) mediopassive form \iwi{CLuw}{\textit{waššāri}} ‘to be good, be dear’ (< \iwi{PIE}{*\textit{h\textsubscript{1}u̯(e)s-ór}} + \textit{-i}) and the HLuw.\ active verb \iwi{HLuw}{\textit{wa/i-sa}-} /was(s)-\textsuperscript{ti}/ ‘id.’.\footnote{For a summary of the research on the PIE “\isi{stative}”, see recently \citealt{Grestenberger2019}. In the case of \iwi{CLuw}{\textit{waššāri}}, it is difficult to reconstruct a \isi{stative} of the type “zero-grade of the root with accent on the suffix” due to the expected effects of Čop’s Law in the root, i.e., *éCV > áCCV. Nonetheless, it is always possible that the shape of the root would have been analogically introduced here after another co-existing stem formation.} A \textit{u}-stem adjective is found in CLuw.\ \iwi{CLuw}{\textit{wāšu}-} ‘good’, HLuw.\ \textit{wa/i-su} (adv.)\iw{HLuw}{\textit{wa/i-su} (adv.)} ‘id.’ and Pal. \iwi{Palaic}{\textit{wāšu}-} ‘id.’, all going back to a generalized \textit{o}-grade of an \isi{acrostatic} paradigm \iwi{PIE}{*\textit{h\textsubscript{1}u̯ós-u-/h\textsubscript{1}u̯és-u}-}. The crucial \textit{nu}-formation is attested in HLuw.\ \iwi{HLuw}{\textit{u-sa-nú-wa/i}-} /usnuwa-\textsuperscript{i}/ ‘to bless’.\footnote{What appears to be a zero-grade might just be a contraction of /wa-/ > /u-/, which sporadically occurs in Iron Age Luwian. Nonetheless, attestations from the earlier stages of the language are required to verify this.} In the end, an apparent “\textit{u}-extended” verb is attested in HLuw.\ \iwi{HLuw}{\textit{wa/i\textsubscript{5}-sú-ha}} ‘I was good, successful’ and probably also in CLuw.\ \iwi{CLuw}{\textit{wašuddu}[} (3sg.imp.act.?).\footnote{The Cuneiform Luwian piece of evidence is found in a highly fragmentary context which calls for caution; see \citealt[179]{Sasseville2021b}.} Unfortunately, due to the lack of attestations, it is not entirely clear what type of conjugation the latter verb has, although, assuming that the Cuneiform Luwian word is complete, it would definitely belong to the non-lenited \isi{\textit{mi}-conjugation} (just like Lyc.\ B \textit{welpu-\textsuperscript{ti}}\iw{Lycian}{\textit{welpu}-}, not **\textit{welpuwe-\textsuperscript{ti}} or **\textit{welpu-\textsuperscript{di}}). The formation in -\textit{nuwa-\textsuperscript{i}}\iw{CLuw}{-\textit{nuwa}-},\iw{HLuw}{-\textit{nuwa}-} despite being primarily analyzed as a deadjectival \is{deadjectival!verb} formation (see \sectref{sec:nu}), may have been reanalyzed\is{reanalysis} at any time as a \isi{causative} formation of the root verb. The next question is what lies behind the stem formation of the “\textit{u}-verb”.\is{*-\textit{u}-present} Sasseville’s explanation as a denominal verb\is{denominal!verb} \is{deadjectival!verb} derived from \iwi{CLuw}{\textit{wāšu}-}\iw{HLuw}{\textit{wāšu}-} ‘good’ falls through (see \sectref{sec:uadj}). A synchronic verbal suffix -\textit{u}- is not possible considering the various conjugations shown by the other data (cf.\ below). What remains is the explanation as a root extension, which would have been triggered by the \isi{reanalysis} of the \isi{causative} formation as an \textit{n}-infixed stem \is{nasal present} based on a root ending in -\textit{u}. Moreover, it is striking that the meaning of the verb is basically the same as that of the two radical verbs, which also points towards a simple \textit{u}-extension, i.e., back-projected \iwi{PIE}{*\textit{h\textsubscript{1}u̯esu}-} or \mbox{\iwi{PIE}{*\textit{h\textsubscript{1}useu̯}-},} backformed\is{backformation} from the \textit{nu}-formation.\is{nasal present} If this reasoning is correct, this example alone would contradict Koch’s “requirement” for the \isi{reanalysis}, because both the adjectival base and the root verb without -\textit{u}- are still present in the Luwian language.

\largerpage
The second root adduced here is \iwi{PIE}{*\textit{h\textsubscript{3}er}-} ‘to arise’: This root has a \isi{stative} verb attested in Hitt.\ \textit{ār-tari}\iw{Hitt}{\textit{ārtari}} ‘he stands’ (< \iwi{PIE}{*\textit{h\textsubscript{3}ér-tor}} + -\textit{i}). There are also a few (de)radical formations found in CLuw.\ \textit{āri-\textsuperscript{ti}}\iw{CLuw}{\textit{āri}-} and Pal. \textit{ariye/a-\textsuperscript{ti}}\iw{Palaic}{\textit{ariye/a}-} ‘to raise’ (< \iwi{PIE}{*\textit{h\textsubscript{3}r̥-i̯é/ó}-}) as well as Hitt.\ \textit{arāi-}/\textit{-i(ya)-\textsuperscript{ḫi}}\iw{Hitt}{\textit{arāi}-} and Lyc.\ \textit{erije-\textsuperscript{ti}}\iw{Lycian}{\textit{erije}-} ‘to arise’ (< \mbox{*\textit{h\textsubscript{3}rói̯h\textsubscript{1}-/}}\textit{h\textsubscript{3}rih\textsubscript{1}-´}\iw{PIE}{*\textit{h\textsubscript{3}rói̯h\textsubscript{1}-/h\textsubscript{3}rih\textsubscript{1}}-´} {<}{<} \mbox{\iwi{PIE}{*\textit{h\textsubscript{3}ór-/h\textsubscript{3}r}-´}).} A \textit{u}-adjective is attested in CLuw.\ and Pal. \textit{arū-/-aw-}\iw{CLuw}{\textit{arū}-}\iw{Palaic}{\textit{arū}-} ‘standing (straight), tall, high (\textit{vel sim}.)’ (< \iwi{PIE}{*\textit{h\textsubscript{3}(é)r-u-/-éu̯}-}).\footnote{The exact semantics of this adjective are not entirely clear, since it may refer to the king as a favorable attribute (KUB 35.165 obv. 23), to crops (KUB 48.69, 1; here the crops are to grow \iwi{Palaic}{\textit{arawa}} in the manner of a man), and to a tree(?) or better to a door (KUB 35.136+ iv 4', 8', 11', 16', 19', 20'; where the CLuw.\ \iwi{CLuw} {\textsuperscript{GIŠ}\textit{ḫīluwa}}, which is modified by \iwi{CLuw}{\textit{arūwa}}, may contextually match Hitt.\ \iwi{Hitt}{\textsuperscript{GIŠ}IG} ‘door’ in line 24').} The well-known \isi{causative} formation in \mbox{\iwi{Hitt}{-\textit{nu}-}} is attested both as Hitt.\ \textit{arnu-\textsuperscript{mi}}\iw{Hitt}{\textit{arnu}-} ‘to set in motion, carry off’ and HLuw.\ \iwi{HLuw}{CRUS+\textit{RA/I-nu-wa/i}-} /arnuwa-\textsuperscript{i}/ ‘(lit.) to let stand, abandon (?)’ and is commonly analyzed as a deverbal formation,\is{deverbal/deverbative!verb} although a deadjectival formation \is{deadjectival!verb} derived from the \textit{u}-adjective is also formally possible.\footnote{The HLuw.\ verb must be derived from \iwi{HLuw}{CRUS+\textit{RA/I}} ‘to stand’, as it is made obvious by the logogram (cf.\ \citetalias{eDiAna}  \href{http://www.ediana.gwi.uni-muenchen.de/dictionary.php?lemma=2317}{ID 2317} on other unpersuasive derivational analyses). The exact meaning of the \isi{causative}/\isi{factitive} verb is debatable in the two contexts, especially since it is used each time with a different \isi{preverb}.} Finally, a “\textit{u}-verb” is attested in CLuw.\ \iwi{CLuw}{\textit{aruwarūwa}-} ‘to lift, erect (?)’, HLuw.\ \iwi{HLuw}{*261.PUGNUS-\textit{ru}-} /aruwaruwa-/ ‘to erect (?)’, both reduplicated, and in Lyc.\ \textit{eruwe-\textsuperscript{ti}}\iw{Lycian}{\textit{eruwe}-}, with a highly debated meaning, but which is interestingly \isi{intransitive}, thus making the Luwian reduplicated formation a \isi{causative} thereof.\footnote{Cf.\ \citetalias{eDiAna}  \href{http://www.ediana.gwi.uni-muenchen.de/dictionary.php?lemma=2231}{ID 2231}.} From an etymological point of view, the Lycian verb should simply mean ‘to stand’ (and the CLuw.\ verb ‘to cause to stand’), although its use with the \isi{preverb} \iwi{Lycian}{\textit{hri}} ‘up’ in a context of funerary instructions and prohibitions leaves space for a more specialized meaning. Here again the “\textit{u}-verb” of the \isi{\textit{ḫi}-conjugation} is morphologically unexplained so far, although the Greek formation \iwi{AG}{ὀρούω} ‘to rise quickly, rush away’ built on the same root (\citetalias[  1107]{Beekes2010}) might correspond to the \textit{u}-verb found in Luwian and Lycian, i.e., \iwi{PIE}{*\textit{h\textsubscript{3}róu̯-/h\textsubscript{3}ru}-´}, with the ablaut pattern of the \isi{\textit{ḫi}-conjugation}.\footnote{R.\ Garnier (pers.\ comm.).} Taken together, this would suggest that the root \iwi{PIE}{*\textit{h\textsubscript{3}er}-} was reinterpreted as \iwi{PIE}{*\textit{h\textsubscript{3}reu̯}-} in both Greek and Luwian with the root vocalism placed between the last consonant and the -\textit{u}-, as in the Avestan example \iwi{Avestan}{\textit{dǝbaoman}-} mentioned above. Following Koch’s reasoning, this extension would have been caused by the reinterpretation of {*\textit{h\textsubscript{3}r̥-néu̯-/-nu}-´} or {*\textit{h\textsubscript{3}r̥-né-u̯-/-n-u}-´}\iw{PIE}{*\textit{h\textsubscript{3}r̥-né(-)u̯-/-n(-)u}-´} (← \iwi{PIE}{*\textit{h\textsubscript{3}(é)r-u-/-éu̯}-}) as an \textit{n}-infixed *\textit{h\textsubscript{3}r̥-n-u̯}- to a verbal root \iwi{PIE}{*\textit{h\textsubscript{3}reu̯}-}.


As a third example, there is the root \iwi{PIE}{*\textit{serk}-} with semantics surrounding the concept of lifting and sublimity found in Hitt.\ \textit{šarkiške/a-\textsuperscript{mi}}\iw{Hitt}{\textit{šarkiške/a}-} ‘to be sublime’ (< \iwi{PIE}{*\textit{sr̥k-s\kinvbreve é/ó}-}), an \isi{imperfective} formation derived from an unattested root verb \iwi{Hitt}{*\textit{šark}-}. Hittite also has a \textit{u}-adjective \textit{šarku-/-aw}-\iw{Hitt}{\textit{šarku}-} ‘sublime’ (< \iwi{PIE}{*\textit{s(é)rk-u-/-éu̯}-}) and a \textit{nu}-formation thereof, i.e., \textit{šarg(a)nu-\textsuperscript{mi}}\iw{Hitt}{\textit{šarg(a)nu}-} ‘to make sublime (?)’ (< \iwi{PIE}{*\textit{sr̥k-né-u̯-/-n-u-}´}), although it is attested in a fragmentary context.\footnote{Thus also \citealt{Luraghi1992}: 156.} Finally, Palaic shows a “\textit{u}-verb” built on the same root, i.e., \iwi{Palaic}{\textit{šarku}-} ‘to lift’, whose conjugation is however obscured by the lack of at least a 3sg. pres. attestation.\footnote{See \citetalias{eDiAna}  \href{http://www.ediana.gwi.uni-muenchen.de/dictionary.php?lemma=2736}{ID 2736}.} As opposed to the last two examples, the “\textit{u}-verb” is \isi{transitive}, whereas the root verb, relying solely on its \isi{imperfective} formation, appears to have been a \isi{stative} verb. It is not impossible that the Palaic verb \iwi{Palaic}{\textit{šarku}-} was originally \isi{intransitive} and secondarily acquired transitivity,\is{transitive} but without further evidence, this claim is unverifiable. Again, it is worth asking how exactly the \textit{u}-verb \is{*-\textit{u}-present} was created and if the -\textit{u}- could have been reanalyzed\is{reanalysis} as part of the root on account of the \textit{nu}-formation, provided that it existed also in Palaic or in its prehistory. Furthermore, the \textit{nu}-formation is, with respect to its derivational base, ambiguous, just like in other examples. Both the \textit{u}-adjective (via infixation) and the unattested root verb (via \isi{suffixation}) are morphologically fine as derivational bases, which again illustrates the constant possibility of morphological \isi{reanalysis}.

More examples can be adduced, even though not all types of formations are attested, which can also be due to chance. The sole example to be discussed here starts with Hitt.\ \textit{nāḫ-/naḫḫ-\textsuperscript{ḫi}}\iw{Hitt}{\textit{nāḫ-/naḫḫ}-} ‘to fear’ and HLuw.\ */nah(h)a-\textsuperscript{i}/ ‘id.’ found as the base of the \isi{imperfective} \iwi{HLuw}{\textit{na-ha-sa}-} (< \iwi{PIE}{*\textit{noh\textsubscript{2}-(h\textsubscript{1})s-}}) going back to the root verb of the \isi{\textit{ḫi}-conjugation} *\textit{nóh\textsubscript{2}-/n\textsubscript{ə}h\textsubscript{2}}-´. Furthermore, it is attested next to the mediopassive Hitt.\ \iwi{Hitt}{\textit{naḫḫantat}} ‘they were an object of fear (?)’.\footnote{\citet[120]{Neu1968} translates it instead as ‘they stood in awe’, thus as showing no semantic distinction to the active root verb. However, since the context is about the capture of Kaskean men, a translation ‘they were an object of fear’ for the \iwi{Hitt}{\textit{naḫḫantat}} is even more probable. However, this cannot be proven beyond doubt, since the context is fragmentary and the restoration of the preceding word [\textit{nu-u}]\textit{n}\textsuperscript{?}-[\textit{n}]\textit{a}\textsuperscript{?}\textit{-aš}\textsuperscript{!} ‘(to) us’ suggested by Neu does not correspond to the readings offered by the autography KBo 50.16.} A \textit{u}-stem occurs as the base of CLuw.\ \iwi{CLuw}{\textit{naḫḫuwašša/i}-} ‘frightening’, which, if it were adjectival, must have been at least substantivized before the \isi{genitival} adjective in \iwi{PIE}{*-\textit{osi̯o}-} was derived (cf.\ CLuw.\ \iwi{CLuw}{\textit{wāšuwašša/i}-} ‘of the goods’). Further evidence for a nominal \textit{u}-stem comes from HLuw.\ \iwi{HLuw}{\textit{na-hu-ti}-} ‘fright’, which belongs to the nominal type in \iwi{HLuw}{*-\textit{u-ti}-} as in Hitt.\ \iwi{Hitt}{-\textit{uzzi}-}. Unfortunately, a formation in \iwi{PAn}{*-\textit{nu}-} is so far missing from the data, but a CLuw.\ “\textit{u}-verb” \textit{naḫḫūwa-\textsuperscript{i}}\iw{CLuw}{\textit{naḫḫūwa}-} ‘to be an object of concern’ going back to a \textit{ḫi}-conjugated\is{\textit{ḫi}-conjugation} *\textit{n\textsubscript{ə}h\textsubscript{2}óu̯-/n\textsubscript{ə}h\textsubscript{2}u-´} (most likely with a restored full grade between \textit{n} and \textit{h\textsubscript{2}}) is attested.\footnote{Unfortunately, it cannot be synchronically proven that the ablaut between \iwi{PIE}{*\textit{-h\textsubscript{2}}-} and \iwi{PIE}{*-\textit{u}-} was marked by means of the recurring plene spellings with the sign <u>, i.e.,   
\includegraphics[]{Glossenkeil}\textit{na-aḫ-ḫu-u-wa}-\iw{HLuw}{\textit{naḫḫūwa}-} … (\citetalias{eDiAna}  \href{http://www.ediana.gwi.uni-muenchen.de/dictionary.php?lemma=3109}{ID 3109}), since the use of this sign following a laryngeal is a strong tendency in Hittite and Luwian (see \citealt[67]{Kimball1999} and \citealt[51--52]{Kloekhorst2008}). Nonetheless, it is worth noticing the same \textit{plene} spelling in \iwi{HLuw}{\textit{a-ru-u-ru-u-wa-am-mi-iš}} (\citetalias{eDiAna}  \href{http://www.ediana.gwi.uni-muenchen.de/dictionary.php?lemma=2232}{ID 2232}), but cf.\ also the adjective \iwi{HLuw}{\textit{a-ru-u}-} ‘high’.} This “\textit{u}-verb”,\is{*-\textit{u}-present} just like the others, cannot be satisfactorily explained as a denominative\is{denominative!verb} via a suffix \iwi{CLuw}{-\textit{u}-} or even via \isi{conversion} from an adjectival \textit{u}-stem, because the \iwi{CLuw}{-\textit{u}-} does not act like a synchronic verbal suffix, which would have to be restricted to a specific conjugation and the same applies \textit{a priori} for the process of \isi{conversion}. The data discussed above, however, demonstrate that these “\textit{u}-verbs”\is{*-\textit{u}-present} can belong either to the \textit{mi}- \is{\textit{mi}-conjugation} or to the \isi{\textit{ḫi}-conjugation}. In his discussion of PIE \textit{u}-presents,\is{*-\textit{u}-present} \citet[141--143]{Jasanoff2003} considered the Anatolian evidence to be inconclusive in this respect, while still pleading for the \isi{\textit{ḫi}-conjugation} as the original one because of the frequent \isi{thematization} of these verbs in Indo-European and of Hitt.\ \textit{lāḫū-\textsuperscript{ḫi}}\iw{Hitt}{\textit{lāḫū}-} ‘to pour’, despite Hitt.\ \mbox{\textit{tarḫu-\textsuperscript{mi}}}\iw{Hitt}{\textit{tarḫu}-}.\footnote{One should take note that a late \textit{nu}-formation built on the verb matching Hitt.\ \textit{tarḫu-\textsuperscript{mi}}\iw{Hitt}{\textit{tarḫu}-} has now surfaced in HLuw.\ \iwi{HLuw}{TONITRUS\textit{-wa/i-nú-wa/i-tu}} /tarhuwanuwa-/ (\citetalias{eDiAna}  \href{http://www.ediana.gwi.uni-muenchen.de/dictionary.php?lemma=2855}{ID 2855}). Here, however, the “\textit{u}-extension” is preserved in the base of the \textit{nu}-formation, showing that the \iwi{HLuw}{-\textit{u}-} is understood to be part of the root in Iron Age Luwian.} In a recent article, \citet[67, fn.\ 24]{Jasanoff2023}  withdraws this view and argues instead that \textit{u}-presents\is{*-\textit{u}-present} were originally proto-middles\is{proto-middle} of the  *\textit{e} : zero-ablauting type, which could secondarily inflect as actives with subsequent \isi{thematization}.\ From his perspective, the actual conjugation that is attested in a Hittite \textit{u}-verb need not reflect the original situation of the PIE “\textit{u}-present”,\is{*-\textit{u}-present} which is seen here as a flawed approach. Adding more examples to the list here, e.g., CLuw.\ \textit{talku-\textsuperscript{ti?}}\iw{CLuw}{\textit{talku}-} ‘to flatten (?)’ (cf.\ Hitt.\ \textit{ištalk-\textsuperscript{mi}}\iw{Hitt}{\textit{ištalk}-} ‘to flatten’), CLuw.\ \textit{malḫu\nobreakdash-\textsuperscript{ti?}}\iw{CLuw}{\textit{malḫu}-} ‘to crush’ (cf.\ Hitt.\ \textit{malla-/mall-\textsuperscript{ḫi}}\iw{Hitt}{\textit{malla-/mall}-} ‘to grind’) and Pal.\ \textit{āp(p)ū-\textsuperscript{i}}\iw{Palaic}{\textit{āp(p)ū}-} ‘?’ (\citetalias{eDiAna}  \href{http://www.ediana.gwi.uni-muenchen.de/dictionary.php?lemma=1566}{ID 1566}), simply confirms the reality of the mixed situation, which should be taken seriously.\footnote{For a deradical formation\is{derivation!deradical} in \iwi{PIE}{*-\textit{i̯e/o}-} \is{*-\textit{i̯e/o}-present}  derived from both the \textit{u}-extended and the bare root, cf.\ Hitt.\ \textit{šarḫiye/a-\textsuperscript{mi}}\iw{Hitt}{\textit{šarḫiye/a}-} ‘to eviscerate, maul’ and Lyc.\ B \textit{zrqqi-\textsuperscript{ti}}\iw{Lycian}{\textit{zrqqi}-} ‘to maul’ going back to \iwi{PIE}{*\textit{sr̥h\textsubscript{3}-i̯é/ó}-} and \iwi{PIE}{*\textit{sr̥h\textsubscript{3}u̯-i̯é/ó}-} respectively (cf.\ also CLuw.\ \iwi{CLuw}{\textit{šarḫuntalla/i}-}, \iwi{CLuw}{\textit{šarḫuwant(i)}-} vs. \iwi{CLuw}{\textit{šarḫadu}}, \mbox{\iwi{CLuw}{\textit{šarḫama/i}-}}, respectively;  \citetalias{eDiAna}  \href{http://www.ediana.gwi.uni-muenchen.de/dictionary.php?lemma=3046}{ID 3046}). It is also noteworthy that \textit{ḫuišuwe/a-\textsuperscript{mi}}\iw{Hitt}{\textit{ḫuišuwe/a}-} ‘to live, stay alive’ (next to \textit{ḫuišu-/-aw-}\iw{Hitt}{\textit{ḫuišu}-} ‘alive, raw’, \textit{ḫuiš-\textsuperscript{mi}}\iw{Hitt}{\textit{ḫuiš}-} ‘to live’ and \textit{ḫuišnu-\textsuperscript{mi}}\iw{Hitt}{\textit{ḫuišnu}-} ‘to keep alive, let live’), which was addressed earlier in this paper, may also reflect a formation in \iwi{PIE}{*-\textit{i̯e/o}-} \is{*-\textit{i̯e/o}-present} derived from a \textit{u}-extended root \iwi{PIE}{*\textit{h\textsubscript{2}u̯esu}-} in the same manner as it happened to \iwi{PIE}{*\textit{h\textsubscript{1}u̯esu}-}. }

The synchronic use of both the \textit{mi}- \is{\textit{mi}-conjugation} and the \isi{\textit{ḫi}-conjugation} is comparable not only to Anatolian root verbs in general, but also to the \textit{n}-infixed verbs \is{nasal present} in Hittite, which may either take the \isi{\textit{mi}-conjugation}, e.g., \textit{ḫarni(n)k-\textsuperscript{mi}}\iw{Hitt}{\textit{ḫarni(n)k}-} ‘to make disappear’, \textit{šarni(n)k-\textsuperscript{mi}}\iw{Hitt}{\textit{šarni(n)k}-} ‘to compensate’, \textit{ištarni(n)k-\textsuperscript{mi}}\iw{Hitt}{\textit{ištarni(n)k}-} ‘to afflict’, or the \isi{\textit{ḫi}-conjugation}, e.g., \textit{šunna-/šunn-\textsuperscript{ḫi}}\iw{Hitt}{\textit{šunna-/šunn}-} ‘to fill’, \textit{tarna-/tarn-\textsuperscript{ḫi}}\iw{Hitt}{\textit{tarna-/tarn}-} ‘to let (go)’.\footnote{For different opinions on whether this situation is adoptable for PIE or whether it is secondary, see recently \citealt[107--108]{Melchert2022} with further literature and \citealt{VillanuevaSvenssonforth}.} Furthermore, as shown in \citealt[481--487]{Sasseville2021}, the two conjugational types can be reconstructed for Proto-Anatolian formations in \iwi{PAn}{*-\textit{nu}-} based on the Anatolian comparative evidence.\footnote{On the problem regarding the conjugation of the suffix \iwi{PIE}{*-\textit{s\kinvbreve e/o}-} in Anatolian, see \citealt[176--177]{Oettinger2023}.} Taken at face value, this situation brings further evidence to the working hypothesis that the \textit{u}-extensions and the \textit{n}-infixations are ultimately related through a process of \isi{reanalysis} involving root verbs, \textit{u}-adjectives and \textit{nu}-formations. Moreover, there would be no need to attempt to find out which conjugation is the original one among “\textit{u}-verbs”, since its determination would be parallel to the one among root verbs, i.e., both conjugational types or, in other words, both types of formations can co-exist in the same root.\footnote{\citet[67, fn.\ 24]{Jasanoff2023} sees the mixed situation attested in Hittite as secondary, analyzing,  e.g., Hitt.\ \textit{lāḫū-\textsuperscript{ḫi}}\iw{Hitt}{\textit{lāḫū}-} ‘to pour’ as an originally \textit{h\textsubscript{2}e}-conjugated \is{*\textit{h\textsubscript{2}e}-conjugation} aorist instead of a \textit{u}-present.\is{*-\textit{u}-present} This approach, however, cannot be synchronically verified. Another idea put forth is that the \isi{\textit{ḫi}-conjugation} of *-\textit{nu}-verbs \is{nasal present} comes from the denominal type\is{denominal!verb} (just like for the \isi{conversion} process for stems in \iwi{PIE}{*-\textit{eh\textsubscript{2}}-}), whereas the deverbal type\is{deverbal/deverbative!verb} was conjugated according to the \isi{\textit{mi}-conjugation} (\citealt[485--486]{Sasseville2021}; \citealt{VillanuevaSvensson2024}, \citeyear{VillanuevaSvenssonforth}). It is worth asking whether one could extend this pattern to root verbs derived from root nouns\is{root!noun} by \isi{conversion}.} 

As it could be deduced from the Anatolian data, the requirement that the \textit{u}-adjective must first disappear in the language for the root to be extended with -\textit{u}- in a given example cannot be maintained. At least the first two examples discussed above clearly demonstrate that Luwian can have the \isi{reanalysis} of the root with \iwi{CLuw}{-\textit{u}-}\iw{HLuw}{-\textit{u}-} while preserving the \textit{u}-adjective. This, however, does not falsify Koch’s hypothesis on the origin of \textit{u}-extensions, but rather indicates that the reason behind it must be sought elsewhere. Looking again at the Anatolian examples shown above, the semantics of the “\textit{u\nobreakdash-}verbs” seem to match the ones of the \isi{stative} (or mediopassive) verbs, i.e., 1) CLuw.\ \iwi{CLuw}{\textit{waššāri}} ‘to be good, be dear’ = HLuw.\ \iwi{HLuw}{\textit{wa/i\textsubscript{5}-sú-ha}} ‘I was good, successful’, 2) Hitt.\ \textit{ār-tari}\iw{Hitt}{\textit{ārtari}} ‘he stands’ = Lyc.\ \textit{eruwe-\textsuperscript{ti}}\iw{Lycian}{\textit{eruwe}-} ‘(intr.)’ with its \isi{causative} reduplicated stem CLuw.\ \iwi{CLuw}{\textit{aruwarūwa}-} ‘to lift, erect (?)’, HLuw.\ \iwi{HLuw}{*261.PUGNUS-\textit{ru}-} /aruwaruwa-/ ‘to erect (?)’ and possibly also 3) Hitt.\ \iwi{Hitt}{\textit{naḫḫantat}} ‘they were an object of fear (?)’ = \textit{naḫḫūwa-\textsuperscript{i}}\iw{HLuw}{\textit{naḫḫūwa}-} ‘to be an object of fear/concern’. This observation implies that there is a strong affinity in these examples between the “\textit{u}-verb” and the \isi{stative} verb, which in turn has a close affinity to the \textit{u}-adjective.\footnote{On this type of relationship  between \textit{u}-adjectives and \isi{stative} verbs in state-oriented verbal roots in Indo-Iranian, see \citealt{Rau2013}.} In the end, the missing link would be the \textit{n}-infix/\textit{nu}-formation.

\section{Conclusion}

In conclusion, this paper has provided  support for the relationship between \textit{u}-adjectives, \textit{nu}-\isi{factitive}s/\isi{causative}s and \textit{u}-extensions by demonstrating that the phenomenon of \isi{reanalysis} is responsible for the creation or extension of certain derivational processes. A difficulty encountered in our data concerns the direction of the \isi{reanalysis} as well as the motivation behind it. To avoid circularity, attention was paid to the rules of word formation to limit the derivational scope. In Anatolian, it could be observed that there is a weak barrier between \textit{nu}-\isi{factitive}s and \textit{nu}-\isi{causative}s.\is{nasal present} A given example can be interpreted either as an \textit{n}-infix derived from a \textit{u}-adjective (i.e., \textit{nu\nobreakdash-}\isi{factitive}) or as a \textit{nu}-suffixed\is{suffixation} formation from a root verb (\textit{nu}-\isi{causative}) and an objective choice is not always possible. Furthermore, it has become clear that the firmly established denominal\is{denominal!verb} \textit{n}-infix \is{nasal present} in the Anatolian languages began to be reanalyzed\is{reanalysis} as a \textit{nu}-suffix and could thus more easily expand its derivational scope.

Building on the PIE “prototype” \iwi{PIE}{*\textit{dʰ(é)bʰ-u-/-éu̯}-} ‘small’ → \iwi{PIE}{*\textit{dʰ(e)bʰ-n-éu̯-/-u}-´} → (root reanalyzed\is{reanalysis} as) \iwi{PIE}{*\textit{dbʰeu̯}-}, the phenomenon of \textit{u}-extended verbs \is{*-\textit{u}-present} in Anatolian was explored with a selected set of data, in order to understand the origin of the stem formation. The data presented can be understood as exemplifying the same derivational chain as our proto-type. As a consequence, the “\textit{u}-verbs” HLuw.\ \iwi{HLuw}{\textit{wa/i\textsubscript{5}-sú-ha}}, Lyc.\ \textit{eruwe-\textsuperscript{ti}}\iw{Lycian}{\textit{eruwe}-} (with CLuw.\ \iwi{CLuw}{\textit{aruwarūwa}-}), Pal.\ \iwi{Palaic}{\textit{šarku}-} and CLuw.\ \textit{naḫḫūwa-\textsuperscript{i}}\iw{CLuw}{\textit{naḫḫūwa}-} were explained as reanalyzed\is{reanalysis} \textit{u}-extended root verbs based on the \isi{reanalysis} of the derivational base of the \textit{nu-/n-}infix formation.\ This Anatolian verbal stem class and its complexity is therefore effectively comparable to the Old Indic 8\textsuperscript{th} class.

\section*{Abbreviations}
\begin{tabularx}{.55\textwidth}{lQ}
CLuw. & Cuneiform Luwian \\
Gk. & Greek \\
Hitt. & Hittite \\
HLuw. & Hieroglyphic Luwian \\
intr. & intransitive\\
KBo & Keilschrifttexte aus Boghazköi \\
KUB & Keilschrifturkunden aus Boghazköi \\
\end{tabularx}
\begin{tabularx}{.43\textwidth}{lQ}
Luw. & Luwian \\
Lyc. & Lycian \\
Lyd. & Lydian \\
OInd. & Old Indic \\
Pal. & Palaic \\
PIE & Proto-Indo-European \\
Ved. & Vedic \\
\end{tabularx}
 
{\sloppy\printbibliography[heading=subbibliography,notkeyword=this]}
\end{document}

 












